\documentclass{article}
\usepackage[margin=0.8in]{geometry}
\usepackage{amsmath,amssymb,amsthm,mathtools}
\usepackage{mathrsfs,bm,bbm}
\usepackage{graphicx,booktabs,array,multirow,tabularx,longtable}
\usepackage{enumitem}
\usepackage{xcolor}
\usepackage{algorithm,algpseudocode}
\usepackage{subcaption,tikz}
\usepackage{siunitx,cancel,accents}
\usepackage{microtype}
\usepackage[numbers,sort&compress]{natbib}
\usepackage[colorlinks=true,linkcolor=blue,citecolor=blue,urlcolor=blue]{hyperref}
\usepackage{cleveref}
\setlength{\emergencystretch}{2em}
\newtheorem{assumption}{Assumption}
\newtheorem{definition}{Definition}
\newtheorem{theorem}{Theorem}
\newtheorem{lemma}{Lemma}
\newtheorem{proposition}{Proposition}
\newtheorem{openendedquestion}{Open-ended Question}
\newtheorem{conjecture}{Conjecture}
\newcommand{\coreref}[1]{\ref{#1}}

\begin{document}

\sloppy

\section*{exp\_\allowbreak{}binary\_\allowbreak{}truthbound\_\allowbreak{}complete}

\noindent\textit{Specialization:} v1\par

\paragraph{TL;DR.} A six-coordinate three-atom experiment has all-observable variance-bound optimum \(107/72\), full Boolean quadratic/cubic optimum \(55/36\), and direct Harshaw--Middleton--Savje positive-semidefinite quadratic optimum \(7/6+\sqrt{17}/9\). Thus the exact quartic gaps are respectively \(1/24\) against every observable Boolean quadratic or cubic rule and \((8\sqrt{17}-23)/72\) against the published quadratic class. The quartic optimizer has an assignmentwise nonnegative rational implementation. More generally, for fixed binary potential outcomes under a finite known randomization, expected observed-data rules are exactly the Boolean Mobius span generated by jointly observable subsets, universal conservativeness is finite pointwise domination, and the observable-margin dual characterizes the Pareto-admissible class. Complement symmetry explains the quartic threshold, and independent replication preserves the block gap. For independent triangular arrays, block independence, bounded scores, and nondegenerate average variance verify the declared Lindeberg--Feller lemma's hypotheses and yield the numerator central limit theorem; the primitive condition \(B^{-2}\sum_j\operatorname{Var}(G_{Bj})\to0\) separately yields the denominator law and asymptotically conservative studentized coverage.

\section{Project justification}

\paragraph{Gap.} The closest arbitrary-design theory, \cite{HarshawMiddletonSavje2026}, characterizes and optimizes quadratic matrix bounds. The zero-pairwise-inclusion correction of \cite{AronowSamii2013} supplies a conservative Horvitz--Thompson rule but not the full observable-function cone or its Pareto frontier; the binary sharp-bound papers \cite{Lu2019FactorialBinary} and \cite{AronowGreenLee2014} treat special designs and different observable information. None supplies the exact higher-order separation or solves this witness's published quadratic semidefinite comparator.

\paragraph{Niche.} Binary support makes every expected observed-data rule a finite multilinear function, so assignment-specific higher-order observability can be characterized without stochastic outcome assumptions and compared exactly with the smaller positive-semidefinite quadratic class.

\paragraph{Fill.} The project gives exact certificates for the all-observable value \(107/72\), the Boolean degree-at-most-three value \(55/36\), and the Harshaw--Middleton--Savje semidefinite value \(7/6+\sqrt{17}/9\). It combines these with the if-and-only-if Boolean support theorem, finite pointwise conservative cone, observable-margin dual, Pareto complete class, rational-certificate corollary, nonnegative implementation, and an exact product law whose degree-restricted part is certified by its own finite linear-program dual. For independent triangular arrays, bounded scores and \(V_B\to V>0\) imply \(B V_B\to\infty\), make every Lindeberg indicator eventually vanish, and therefore permit direct invocation of the declared independent-array Lindeberg--Feller lemma. The primitive row-average variance condition on the realized bound tables then yields the denominator weak law by Chebyshev's inequality and licenses conservative Wald coverage.

\section{Related work}

\cite{HarshawMiddletonSavje2026}, especially Proposition 1, Theorems 1 and 3, Definitions 1--2, and Section 8, supplies the quadratic characterization and supporting-objective theory and explicitly restricts attention to quadratic bounds. Its pairwise design-compatibility criterion embeds every published candidate into the Boolean quadratic cone, while its positive-semidefinite domination requirement makes that class strictly smaller than cube-wise Boolean domination. For the six-coordinate witness, the present exact rank-one primal and algebraic dual-slack certificate solves the published semidefinite program itself at \(7/6+\sqrt{17}/9\), not merely at the containing Boolean lower bound \(55/36\). The resulting direct gap from the quartic optimum is \((8\sqrt{17}-23)/72\). The Harshaw--Middleton--Savje use of \cite{PaluckShepherdAronow2016} is a quadratic school-data empirical illustration, not a reanalysis of the original causal conclusions. \cite{Middleton2021Unifying}, \emph{Unifying Design-based Inference: On Bounding and Estimating the Variance of any Linear Estimator in any Experimental Design}, supplies the arbitrary design--estimator--bound triplet but remains matrix-quadratic. Proposition 2 and Corollary 2 of \cite{AronowSamii2013} use Young's inequality to construct a weakly conservative Horvitz--Thompson correction when pairwise inclusion probabilities vanish; its binary restriction lies in \(\mathcal C_{\mathcal J}(f)\cap\mathcal E_{\mathcal J,2}\), but that paper does not characterize the full Boolean observable-margin cone, its Pareto complete class, or the strict quartic separation. \cite{Lu2019FactorialBinary} supplies a sharp binary inclusion--exclusion correction for the \(2^2\) factorial design, while \cite{AronowGreenLee2014} supplies sharp marginal variance bounds for difference-in-means under complete randomization. \cite{ChattopadhyayImbens2025} develops imputation and contrast variance estimators, including nonquadratic procedures, but does not characterize the complete observable-function space or its Pareto frontier. \cite{AronowSamii2017Interference} fixes the design--exposure--estimand semantics and supplies inverse-probability variance estimation and local-dependence asymptotics for the school-network consumer. \cite{LiDing2017FPCLT} supplies the finite-population limit-theory benchmark. The supporting inference theorem here instead derives the Lindeberg condition from bounded scores and (B V\_B\to\infty), then invokes the cited independent triangular-array Lindeberg--Feller theorem; it derives the denominator weak law by Chebyshev's inequality from \(B^{-2}\sum_j\operatorname{Var}(G_{Bj})\to0\), with blockwise conservativeness kept as a distinct premise. \cite{AzrielKriegerKapelner2026}, \emph{Optimal Designs with Robust Inference for Binary Treatment Effects}, studies balanced designs and conservative CMH- and Robins-type inference under random binary potential outcomes, not unrestricted fixed binary schedules under a fixed arbitrary design.

\section{Setup}

Target (estimand / structural parameter): \(\tau(\theta)=c^\top\theta\).
Primary functional / estimator / diagnostic: \(F_{\mathcal J}(f,q)=\min_{b\in\mathcal E_{\mathcal J}:b\ge f}\mathbb E_q[b]=\max_{\mu:\mu|_{O_z}=q|_{O_z}\;\forall z}\mathbb E_\mu[f]\).
\begin{itemize}
  \item \texttt{K} --- \texttt{coordinate\_\allowbreak{}count}; \(\mathbb N_0\); \texttt{finite\_\allowbreak{}dimension}
  \par\smallskip\noindent\emph{Definition.} \texttt{Number of binary potential-\allowbreak{}outcome coordinates.}
  \item \texttt{I} --- \texttt{coordinate\_\allowbreak{}set}; \(2^{\mathbb N_0}\); \texttt{index\_\allowbreak{}set}
  \par\smallskip\noindent\emph{Definition.} Coordinate set \(I=\{0,\ldots,K-1\}\).
  \item \texttt{Z} --- \texttt{assignment\_\allowbreak{}support}; \(2^{\mathbb N}\); \texttt{design\_\allowbreak{}support}
  \par\smallskip\noindent\emph{Definition.} \texttt{Nonempty finite support of the randomization law.}
  \item \texttt{z} --- \texttt{assignment\_\allowbreak{}atom}; \(Z\); \texttt{assignment\_\allowbreak{}index}
  \par\smallskip\noindent\emph{Definition.} \texttt{Generic assignment atom.}
  \item \texttt{p} --- \texttt{assignment\_\allowbreak{}probability\_\allowbreak{}vector}; \(\Delta(Z)\); \texttt{design\_\allowbreak{}law}
  \par\smallskip\noindent\emph{Definition.} Known randomization law \(p=(p_z)_{z\in Z}\), with \(Z\) defined as its support.
  \item \texttt{O\_\allowbreak{}z} --- \texttt{observed\_\allowbreak{}coordinate\_\allowbreak{}set}; \(2^I\); \texttt{observation\_\allowbreak{}map}
  \par\smallskip\noindent\emph{Definition.} Coordinates revealed under assignment \(z\).
  \item \texttt{v\_\allowbreak{}z} --- \texttt{score\_\allowbreak{}vector}; \(\mathbb R^K\); \texttt{estimator\_\allowbreak{}score}
  \par\smallskip\noindent\emph{Definition.} Known coefficient vector used by the linear effect estimator under assignment \(z\).
  \item \texttt{\textbackslash{}\allowbreak{}Theta} --- \texttt{schedule\_\allowbreak{}space}; \(\{0,1\}^I\); \texttt{fixed\_\allowbreak{}finite\_\allowbreak{}population}
  \par\smallskip\noindent\emph{Definition.} \texttt{Finite binary potential-\allowbreak{}outcome schedule space.}
  \item \texttt{\textbackslash{}\allowbreak{}theta} --- \texttt{potential\_\allowbreak{}outcome\_\allowbreak{}schedule}; \(\Theta\); \texttt{causal\_\allowbreak{}schedule}
  \par\smallskip\noindent\emph{Definition.} \texttt{Fixed schedule;\allowbreak{} it is not sampled from a superpopulation.}
  \item \texttt{S} --- \texttt{coordinate\_\allowbreak{}subset}; \(2^I\); \texttt{mobius\_\allowbreak{}index}
  \par\smallskip\noindent\emph{Definition.} \texttt{Generic subset of coordinates.}
  \item \texttt{m\_\allowbreak{}S} --- \texttt{boolean\_\allowbreak{}monomial}; \(\Theta\to\{0,1\}\); \texttt{boolean\_\allowbreak{}basis}
  \par\smallskip\noindent\emph{Definition.} Squarefree monomial \(m_S(\theta)=\prod_{i\in S}\theta_i\), with \(m_\varnothing=1\).
  \item \texttt{\textbackslash{}\allowbreak{}mathcal J} --- \texttt{observable\_\allowbreak{}complex}; \(2^{2^I}\); \texttt{identified\_\allowbreak{}supports}
  \par\smallskip\noindent\emph{Definition.} \texttt{Downward-\allowbreak{}closed family of subsets observable together under some assignment.}
  \item \texttt{\textbackslash{}\allowbreak{}mathcal E\_\allowbreak{}\{\textbackslash{}\allowbreak{}mathcal J\}\allowbreak{}} --- \texttt{observable\_\allowbreak{}function\_\allowbreak{}space}; \(\mathbb R^\Theta\); \texttt{estimable\_\allowbreak{}expectations}
  \par\smallskip\noindent\emph{Definition.} \texttt{Linear span of observable Boolean monomials.}
  \item \texttt{r} --- \texttt{degree\_\allowbreak{}limit}; \(\mathbb N_0\); \texttt{degree\_\allowbreak{}filtration}
  \par\smallskip\noindent\emph{Definition.} \texttt{Maximum Boolean polynomial degree.}
  \item \texttt{\textbackslash{}\allowbreak{}mathcal E\_\allowbreak{}\{\textbackslash{}\allowbreak{}mathcal J,\allowbreak{}r\}\allowbreak{}} --- \texttt{degree\_\allowbreak{}restricted\_\allowbreak{}observable\_\allowbreak{}space}; \(\mathbb R^\Theta\); \texttt{quadratic\_\allowbreak{}cubic\_\allowbreak{}comparator}
  \par\smallskip\noindent\emph{Definition.} Subspace generated by observable monomials of degree at most \(r\).
  \item \texttt{X\_\allowbreak{}z} --- \texttt{linear\_\allowbreak{}score}; \(\mathbb R\); \texttt{effect\_\allowbreak{}estimator}
  \par\smallskip\noindent\emph{Definition.} Realized score \(X_z(\theta)=v_z^\top\theta\).
  \item \texttt{c} --- \texttt{target\_\allowbreak{}coefficient\_\allowbreak{}vector}; \(\mathbb R^K\); \texttt{estimand\_\allowbreak{}weight}
  \par\smallskip\noindent\emph{Definition.} Design mean score coefficient \(c=\sum_z p_zv_z\).
  \item \texttt{\textbackslash{}\allowbreak{}tau} --- \texttt{finite\_\allowbreak{}population\_\allowbreak{}causal\_\allowbreak{}contrast}; \(\Theta\to\mathbb R\); \texttt{causal\_\allowbreak{}estimand}
  \par\smallskip\noindent\emph{Definition.} Target \(\tau(\theta)=c^\top\theta\).
  \item \texttt{A} --- \texttt{variance\_\allowbreak{}matrix}; \(\mathbb R^{K\times K}\); \texttt{true\_\allowbreak{}variance\_\allowbreak{}coefficients}
  \par\smallskip\noindent\emph{Definition.} Matrix \(A=\sum_zp_zv_zv_z^\top-cc^\top\).
  \item \texttt{f} --- \texttt{design\_\allowbreak{}variance\_\allowbreak{}function}; \(\Theta\to\mathbb R_+\); \texttt{bound\_\allowbreak{}target}
  \par\smallskip\noindent\emph{Definition.} True variance \(f(\theta)=\theta^\top A\theta=\operatorname{Var}_p\{X_z(\theta)\}\).
  \item \texttt{H} --- \texttt{quadratic\_\allowbreak{}bound\_\allowbreak{}matrix}; \(\mathbb R^{K\times K}\); \texttt{published\_\allowbreak{}quadratic\_\allowbreak{}comparator}
  \par\smallskip\noindent\emph{Definition.} \texttt{Generic symmetric matrix used to represent a Harshaw-\allowbreak{}-\allowbreak{}Middleton-\allowbreak{}-\allowbreak{}Savje design-\allowbreak{}compatible quadratic candidate.}
  \item \texttt{h\_\allowbreak{}H} --- \texttt{restricted\_\allowbreak{}quadratic\_\allowbreak{}function}; \(\Theta\to\mathbb R\); \texttt{published\_\allowbreak{}quadratic\_\allowbreak{}comparator}
  \par\smallskip\noindent\emph{Definition.} Binary-cube restriction \(h_H(\theta)=\theta^\top H\theta\).
  \item \texttt{g\_\allowbreak{}z\textasciicircum{}\{\textbackslash{}\allowbreak{}mathrm\{AS\}\allowbreak{}\}\allowbreak{}} --- \texttt{aronow\_\allowbreak{}samii\_\allowbreak{}assignment\_\allowbreak{}table}; \(\{0,1\}^{O_z}\to\mathbb R\); \texttt{published\_\allowbreak{}zero\_\allowbreak{}observability\_\allowbreak{}comparator}
  \par\smallskip\noindent\emph{Definition.} \texttt{Assignment table obtained by restricting the simplified Aronow-\allowbreak{}-\allowbreak{}Samii zero-\allowbreak{}pairwise-\allowbreak{}inclusion correction to binary coordinates.}
  \item \texttt{b\_\allowbreak{}\{\textbackslash{}\allowbreak{}mathrm\{AS\}\allowbreak{}\}\allowbreak{}} --- \texttt{aronow\_\allowbreak{}samii\_\allowbreak{}expected\_\allowbreak{}correction}; \(\Theta\to\mathbb R\); \texttt{published\_\allowbreak{}zero\_\allowbreak{}observability\_\allowbreak{}comparator}
  \par\smallskip\noindent\emph{Definition.} Design expectation \(b_{\mathrm{AS}}(\theta)=\sum_zp_zg_z^{\mathrm{AS}}(\theta|_{O_z})\).
  \item \texttt{g\_\allowbreak{}z} --- \texttt{assignment\_\allowbreak{}truth\_\allowbreak{}table}; \(\{0,1\}^{O_z}\to\mathbb R\); \texttt{variance\_\allowbreak{}bound\_\allowbreak{}estimator}
  \par\smallskip\noindent\emph{Definition.} Observed-data table used when assignment \(z\) occurs.
  \item \texttt{b\_\allowbreak{}g} --- \texttt{expected\_\allowbreak{}bound\_\allowbreak{}function}; \(\Theta\to\mathbb R\); \texttt{expected\_\allowbreak{}variance\_\allowbreak{}bound}
  \par\smallskip\noindent\emph{Definition.} Design expectation \(b_g(\theta)=\sum_zp_zg_z(\theta|_{O_z})\).
  \item \texttt{b} --- \texttt{candidate\_\allowbreak{}bound}; \(\mathcal E_{\mathcal J}\); \texttt{primal\_\allowbreak{}variable}
  \par\smallskip\noindent\emph{Definition.} \texttt{Generic observable expected-\allowbreak{}bound function.}
  \item \texttt{b'} --- \texttt{candidate\_\allowbreak{}bound}; \(\mathcal E_{\mathcal J}\); \texttt{domination\_\allowbreak{}comparator}
  \par\smallskip\noindent\emph{Definition.} \texttt{Generic competing observable expected-\allowbreak{}bound function.}
  \item \texttt{q} --- \texttt{schedule\_\allowbreak{}weighting\_\allowbreak{}distribution}; \(\Delta^\circ(\Theta)\); \texttt{supporting\_\allowbreak{}objective}
  \par\smallskip\noindent\emph{Definition.} Prespecified full-support objective weights; \(q\) is not a law for the fixed schedule.
  \item \texttt{\textbackslash{}\allowbreak{}mathcal C\_\allowbreak{}\{\textbackslash{}\allowbreak{}mathcal J\}\allowbreak{}(f)} --- \texttt{conservative\_\allowbreak{}observable\_\allowbreak{}cone}; \(2^{\mathcal E_{\mathcal J}}\); \texttt{feasible\_\allowbreak{}bounds}
  \par\smallskip\noindent\emph{Definition.} Observable functions that dominate \(f\) on the binary cube.
  \item \texttt{F\_\allowbreak{}\{\textbackslash{}\allowbreak{}mathcal J\}\allowbreak{}} --- \texttt{optimal\_\allowbreak{}bound\_\allowbreak{}value}; \(\mathbb R_+\); \texttt{all\_\allowbreak{}degree\_\allowbreak{}optimum}
  \par\smallskip\noindent\emph{Definition.} Minimum \(q\)-average observable conservative bound.
  \item \texttt{F\_\allowbreak{}\{\textbackslash{}\allowbreak{}mathcal J,\allowbreak{}r\}\allowbreak{}} --- \texttt{restricted\_\allowbreak{}optimal\_\allowbreak{}bound\_\allowbreak{}value}; \(\mathbb R_+\); \texttt{degree\_\allowbreak{}frontier}
  \par\smallskip\noindent\emph{Definition.} Minimum with degree restricted to \(r\).
  \item \texttt{\textbackslash{}\allowbreak{}beta\_\allowbreak{}S} --- \texttt{mobius\_\allowbreak{}coefficient}; \(\mathbb R\); \texttt{observable\_\allowbreak{}support\_\allowbreak{}coefficient}
  \par\smallskip\noindent\emph{Definition.} Coefficient of \(m_S\) in the unique multilinear expansion of a function on \(\Theta\).
  \item \texttt{\textbackslash{}\allowbreak{}pi\_\allowbreak{}S} --- \texttt{joint\_\allowbreak{}observation\_\allowbreak{}probability}; \(\mathbb R_+\); \texttt{implementation\_\allowbreak{}weight}
  \par\smallskip\noindent\emph{Definition.} Probability \(\pi_S=\sum_{z:S\subseteq O_z}p_z\).
  \item \texttt{\textbackslash{}\allowbreak{}mu} --- \texttt{dual\_\allowbreak{}schedule\_\allowbreak{}distribution}; \(\mathbb R_+^\Theta\); \texttt{dual\_\allowbreak{}certificate}
  \par\smallskip\noindent\emph{Definition.} \texttt{Nonnegative weights on schedules used in the observable-\allowbreak{}margin dual.}
  \item \texttt{\textbackslash{}\allowbreak{}mathcal D\_\allowbreak{}\{\textbackslash{}\allowbreak{}mathcal J\}\allowbreak{}(q)} --- \texttt{dual\_\allowbreak{}margin\_\allowbreak{}polytope}; \(2^{\mathbb R_+^\Theta}\); \texttt{dual\_\allowbreak{}feasible\_\allowbreak{}set}
  \par\smallskip\noindent\emph{Definition.} Schedule weights matching \(q\) on every observable Boolean moment.
  \item \texttt{i} --- \texttt{coordinate\_\allowbreak{}index}; \(I\); \texttt{index}
  \par\smallskip\noindent\emph{Definition.} \texttt{Generic coordinate index.}
  \item \texttt{k} --- \texttt{coordinate\_\allowbreak{}index}; \(I\); \texttt{index}
  \par\smallskip\noindent\emph{Definition.} \texttt{Second generic coordinate index.}
  \item \texttt{T} --- \texttt{coordinate\_\allowbreak{}subset}; \(2^I\); \texttt{inversion\_\allowbreak{}index}
  \par\smallskip\noindent\emph{Definition.} \texttt{Generic subset used in Mobius inversion.}
  \item \texttt{\textbackslash{}\allowbreak{}mathbf 1\_\allowbreak{}T} --- \texttt{vertex\_\allowbreak{}schedule}; \(\Theta\); \texttt{mobius\_\allowbreak{}vertex}
  \par\smallskip\noindent\emph{Definition.} Binary schedule with ones exactly on \(T\).
  \item \texttt{d} --- \texttt{parity\_\allowbreak{}index}; \(\mathbb N_0\); \texttt{degree\_\allowbreak{}index}
  \par\smallskip\noindent\emph{Definition.} \texttt{Half-\allowbreak{}degree index in the complement-\allowbreak{}symmetry theorem.}
  \item \texttt{\textbackslash{}\allowbreak{}mathbf 1\_\allowbreak{}K} --- \texttt{all\_\allowbreak{}one\_\allowbreak{}schedule}; \(\Theta\); \texttt{complement\_\allowbreak{}map}
  \par\smallskip\noindent\emph{Definition.} \texttt{The all-\allowbreak{}one binary schedule.}
  \item \texttt{x\_\allowbreak{}i} --- \texttt{centered\_\allowbreak{}binary\_\allowbreak{}coordinate}; \(\{-1,1\}\); \texttt{parity\_\allowbreak{}basis}
  \par\smallskip\noindent\emph{Definition.} Centered coordinate \(x_i=2\theta_i-1\).
  \item \texttt{x} --- \texttt{centered\_\allowbreak{}schedule}; \(\{-1,1\}^6\); \texttt{parity\_\allowbreak{}argument}
  \par\smallskip\noindent\emph{Definition.} Vector collecting the centered coordinates \(x_i\) in the witness.
  \item \texttt{P} --- \texttt{witness\_\allowbreak{}quadratic\_\allowbreak{}polynomial}; \(\{-1,1\}^6\to\mathbb Z\); \texttt{certificate\_\allowbreak{}component}
  \par\smallskip\noindent\emph{Definition.} \texttt{Quadratic centered-\allowbreak{}coordinate polynomial in the witness certificates.}
  \item \texttt{\textbackslash{}\allowbreak{}mathfrak E\_\allowbreak{}\textbackslash{}\allowbreak{}star} --- \texttt{six\_\allowbreak{}coordinate\_\allowbreak{}experiment}; \texttt{strict\_\allowbreak{}separation\_\allowbreak{}witness}
  \par\smallskip\noindent\emph{Definition.} \texttt{The specified three-\allowbreak{}atom witness experiment.}
  \item \texttt{b\textasciicircum{}\textbackslash{}\allowbreak{}star} --- \texttt{quartic\_\allowbreak{}bound}; \(\mathcal E_{\mathcal J}\); \texttt{quartic\_\allowbreak{}primal\_\allowbreak{}certificate}
  \par\smallskip\noindent\emph{Definition.} All-degree optimizer in \(\mathfrak E_\star\).
  \item \texttt{r\textasciicircum{}\textbackslash{}\allowbreak{}star} --- \texttt{quadratic\_\allowbreak{}bound}; \(\mathcal E_{\mathcal J,2}\); \texttt{quadratic\_\allowbreak{}primal\_\allowbreak{}certificate}
  \par\smallskip\noindent\emph{Definition.} Degree-two optimizer in \(\mathfrak E_\star\).
  \item \texttt{h\_\allowbreak{}z\textasciicircum{}\textbackslash{}\allowbreak{}star} --- \texttt{nonnegative\_\allowbreak{}assignment\_\allowbreak{}table}; \(\{0,1\}^{O_z}\to\mathbb Q_+\); \texttt{implementability\_\allowbreak{}certificate}
  \par\smallskip\noindent\emph{Definition.} Assignment-level implementation of \(b^\star\).
  \item \texttt{\textbackslash{}\allowbreak{}mu\textasciicircum{}\textbackslash{}\allowbreak{}star} --- \texttt{all\_\allowbreak{}degree\_\allowbreak{}dual\_\allowbreak{}distribution}; \(\mathbb Q_+^\Theta\); \texttt{rational\_\allowbreak{}dual\_\allowbreak{}certificate}
  \par\smallskip\noindent\emph{Definition.} \texttt{Uniform distribution on the sixteen schedules in the all-\allowbreak{}degree dual certificate.}
  \item \texttt{\textbackslash{}\allowbreak{}nu\textasciicircum{}\textbackslash{}\allowbreak{}star} --- \texttt{degree\_\allowbreak{}three\_\allowbreak{}dual\_\allowbreak{}distribution}; \(\mathbb Q_+^\Theta\); \texttt{rational\_\allowbreak{}dual\_\allowbreak{}certificate}
  \par\smallskip\noindent\emph{Definition.} \texttt{Uniform distribution on the eight schedules in the degree-\allowbreak{}three dual certificate.}
  \item \texttt{B} --- \texttt{block\_\allowbreak{}count}; \(\mathbb N\); \texttt{replication\_\allowbreak{}index}
  \par\smallskip\noindent\emph{Definition.} \texttt{Number of independently randomized component experiments.}
  \item \texttt{j} --- \texttt{block\_\allowbreak{}index}; \(\{1,\ldots,B\}\); \texttt{index}
  \par\smallskip\noindent\emph{Definition.} \texttt{Generic block index.}
  \item \texttt{a\_\allowbreak{}j} --- \texttt{deterministic\_\allowbreak{}block\_\allowbreak{}weight}; \(\mathbb R\); \texttt{aggregation\_\allowbreak{}weight}
  \par\smallskip\noindent\emph{Definition.} Known coefficient multiplying block \(j\)'s score.
  \item \texttt{X\_\allowbreak{}j} --- \texttt{component\_\allowbreak{}score}; \(\mathbb R\); \texttt{product\_\allowbreak{}score}
  \par\smallskip\noindent\emph{Definition.} Linear score of component experiment \(j\).
  \item \texttt{F\_\allowbreak{}j} --- \texttt{component\_\allowbreak{}optimum}; \(\mathbb R_+\); \texttt{product\_\allowbreak{}component\_\allowbreak{}value}
  \par\smallskip\noindent\emph{Definition.} All-degree optimum in component \(j\).
  \item \texttt{F\_\allowbreak{}\{j,\allowbreak{}r\}\allowbreak{}} --- \texttt{restricted\_\allowbreak{}component\_\allowbreak{}optimum}; \(\mathbb R_+\); \texttt{product\_\allowbreak{}component\_\allowbreak{}value}
  \par\smallskip\noindent\emph{Definition.} Degree \(r\) optimum in component \(j\).
  \item \texttt{F\_\allowbreak{}\{\textbackslash{}\allowbreak{}mathrm\{prod\}\allowbreak{}\}\allowbreak{}} --- \texttt{product\_\allowbreak{}optimum}; \(\mathbb R_+\); \texttt{product\_\allowbreak{}value}
  \par\smallskip\noindent\emph{Definition.} \texttt{All-\allowbreak{}degree optimum for the independent product experiment.}
  \item \texttt{F\_\allowbreak{}\{\textbackslash{}\allowbreak{}mathrm\{prod\}\allowbreak{},\allowbreak{}r\}\allowbreak{}} --- \texttt{restricted\_\allowbreak{}product\_\allowbreak{}optimum}; \(\mathbb R_+\); \texttt{product\_\allowbreak{}value}
  \par\smallskip\noindent\emph{Definition.} Degree \(r\) optimum for the independent product experiment.
  \item \texttt{X\_\allowbreak{}\{Bj\}\allowbreak{}} --- \texttt{block\_\allowbreak{}effect\_\allowbreak{}score}; \(\mathbb R\); \texttt{triangular\_\allowbreak{}array\_\allowbreak{}numerator}
  \par\smallskip\noindent\emph{Definition.} Assignment-random block score in row \(B\), with fixed potential outcomes.
  \item \texttt{G\_\allowbreak{}\{Bj\}\allowbreak{}} --- \texttt{block\_\allowbreak{}bound\_\allowbreak{}estimator}; \(\mathbb R\); \texttt{triangular\_\allowbreak{}array\_\allowbreak{}denominator}
  \par\smallskip\noindent\emph{Definition.} \texttt{Assignment-\allowbreak{}level estimator of a conservative expected block variance bound.}
  \item \texttt{\textbackslash{}\allowbreak{}tau\_\allowbreak{}\{Bj\}\allowbreak{}} --- \texttt{block\_\allowbreak{}target}; \(\mathbb R\); \texttt{block\_\allowbreak{}estimand}
  \par\smallskip\noindent\emph{Definition.} Design expectation \(\tau_{Bj}=\mathbb E X_{Bj}\).
  \item \texttt{\textbackslash{}\allowbreak{}sigma\_\allowbreak{}\{Bj\}\allowbreak{}\textasciicircum{}2} --- \texttt{block\_\allowbreak{}variance}; \(\mathbb R_+\); \texttt{true\_\allowbreak{}block\_\allowbreak{}variance}
  \par\smallskip\noindent\emph{Definition.} Design variance \(\sigma_{Bj}^2=\operatorname{Var}(X_{Bj})\).
  \item \texttt{b\_\allowbreak{}\{Bj\}\allowbreak{}} --- \texttt{expected\_\allowbreak{}block\_\allowbreak{}bound}; \(\mathbb R_+\); \texttt{conservative\_\allowbreak{}block\_\allowbreak{}target}
  \par\smallskip\noindent\emph{Definition.} Design mean \(b_{Bj}=\mathbb E G_{Bj}\).
  \item \texttt{V\_\allowbreak{}B} --- \texttt{average\_\allowbreak{}true\_\allowbreak{}variance}; \(\mathbb R_+\); \texttt{numerator\_\allowbreak{}scale}
  \par\smallskip\noindent\emph{Definition.} \(V_B=B^{-1}\sum_j\sigma_{Bj}^2\).
  \item \texttt{W\_\allowbreak{}B} --- \texttt{average\_\allowbreak{}expected\_\allowbreak{}bound}; \(\mathbb R_+\); \texttt{denominator\_\allowbreak{}limit}
  \par\smallskip\noindent\emph{Definition.} \(W_B=B^{-1}\sum_jb_{Bj}\).
  \item \texttt{\textbackslash{}\allowbreak{}widehat W\_\allowbreak{}B} --- \texttt{average\_\allowbreak{}realized\_\allowbreak{}bound}; \(\mathbb R\); \texttt{variance\_\allowbreak{}estimator\_\allowbreak{}scale}
  \par\smallskip\noindent\emph{Definition.} \(\widehat W_B=B^{-1}\sum_jG_{Bj}\).
  \item \texttt{\textbackslash{}\allowbreak{}widehat\textbackslash{}\allowbreak{}tau\_\allowbreak{}B} --- \texttt{average\_\allowbreak{}effect\_\allowbreak{}estimator}; \(\mathbb R\); \texttt{studentized\_\allowbreak{}estimator}
  \par\smallskip\noindent\emph{Definition.} \(\widehat\tau_B=B^{-1}\sum_jX_{Bj}\).
  \item \texttt{\textbackslash{}\allowbreak{}bar\textbackslash{}\allowbreak{}tau\_\allowbreak{}B} --- \texttt{average\_\allowbreak{}causal\_\allowbreak{}target}; \(\mathbb R\); \texttt{triangular\_\allowbreak{}array\_\allowbreak{}estimand}
  \par\smallskip\noindent\emph{Definition.} \(\bar\tau_B=B^{-1}\sum_j\tau_{Bj}\).
  \item \texttt{M} --- \texttt{finite\_\allowbreak{}score\_\allowbreak{}bound}; \(\mathbb R_+\); \texttt{lindeberg\_\allowbreak{}constant}
  \par\smallskip\noindent\emph{Definition.} \texttt{Uniform deterministic bound on block scores.}
  \item \texttt{V} --- \texttt{positive\_\allowbreak{}variance\_\allowbreak{}limit}; \(\mathbb R_{>0}\); \texttt{asymptotic\_\allowbreak{}variance}
  \par\smallskip\noindent\emph{Definition.} Positive limit of \(V_B\).
  \item \texttt{W} --- \texttt{finite\_\allowbreak{}bound\_\allowbreak{}limit}; \(\mathbb R_+\); \texttt{studentization\_\allowbreak{}limit}
  \par\smallskip\noindent\emph{Definition.} Finite limit of \(W_B\), when it exists.
  \item \texttt{\textbackslash{}\allowbreak{}alpha} --- \texttt{significance\_\allowbreak{}level}; \((0,1)\); \texttt{coverage\_\allowbreak{}level}
  \par\smallskip\noindent\emph{Definition.} \texttt{Nominal error probability.}
  \item \texttt{\textbackslash{}\allowbreak{}Phi} --- \texttt{standard\_\allowbreak{}normal\_\allowbreak{}cdf}; \(\mathbb R\to(0,1)\); \texttt{limit\_\allowbreak{}cdf}
  \par\smallskip\noindent\emph{Definition.} \texttt{Standard normal distribution function.}
  \item \texttt{z\_\allowbreak{}\textbackslash{}\allowbreak{}alpha} --- \texttt{normal\_\allowbreak{}critical\_\allowbreak{}value}; \(\mathbb R_+\); \texttt{wald\_\allowbreak{}cutoff}
  \par\smallskip\noindent\emph{Definition.} \(z_\alpha=\Phi^{-1}(1-\alpha/2)\).
  \item \texttt{I\_\allowbreak{}B} --- \texttt{wald\_\allowbreak{}interval}; \(2^{\mathbb R}\); \texttt{confidence\_\allowbreak{}interval}
  \par\smallskip\noindent\emph{Definition.} Conservative studentized interval in row \(B\).
  \item \texttt{\textbackslash{}\allowbreak{}mathcal H\_\allowbreak{}\{\textbackslash{}\allowbreak{}mathrm\{school\}\allowbreak{}\}\allowbreak{}} --- \texttt{decomposition\_\allowbreak{}handle}; \texttt{open\_\allowbreak{}application\_\allowbreak{}handle}
  \par\smallskip\noindent\emph{Definition.} \texttt{Orbit-\allowbreak{}and-\allowbreak{}separator route for an exact school-\allowbreak{}network observable-\allowbreak{}margin program.}
\end{itemize}

\section{Assumptions}

\begin{assumption}[ass:observable-score]\label{ass:observable-score}
\(\operatorname{supp}(v_z)\subseteq O_z\) for every \(z\in Z\).
\par\smallskip\noindent\emph{Standard} (design compatibility of the linear score, \cite{HarshawMiddletonSavje2026}).
\end{assumption}

\begin{assumption}[ass:complement-invariant-variance]\label{ass:complement-invariant-variance}
\(A\mathbf 1_K=0\).
\par\smallskip\noindent\emph{Novel.} This proposal-specific outcome-complement symmetry has a direct design interpretation: it is satisfied when every assignment-specific score has the same total coordinate loading, including zero-sum contrast scores, because complementing all binary outcomes reflects every realized score around the same constant and leaves its randomization variance unchanged.
\end{assumption}

\begin{assumption}[ass:complement-invariant-objective]\label{ass:complement-invariant-objective}
\(q_\theta=q_{\mathbf 1_K-\theta}\) for every \(\theta\in\Theta\).
\par\smallskip\noindent\emph{Novel.} Uniform weighting and any paired prior weighting on binary schedules satisfy this condition; it formalizes outcome-label symmetry without restricting the fixed realized schedule.
\end{assumption}

\begin{assumption}[ass:block-independence]\label{ass:block-independence}
The pairs \(\{(X_{Bj},G_{Bj}):1\le j\le B\}\) are independent across \(j\) under the row \(B\) randomization law.
\par\smallskip\noindent\emph{Standard} (independent randomized blocks, \cite{LiDing2017FPCLT}).
\end{assumption}

\begin{assumption}[ass:uniform-score-bound]\label{ass:uniform-score-bound}
\(|X_{Bj}|\le M<\infty\) for every \(B\), \(j\), and assignment atom.
\par\smallskip\noindent\emph{Standard} (uniform boundedness implying the Lindeberg condition, \cite{LiDing2017FPCLT}).
\end{assumption}

\begin{assumption}[ass:nondegenerate-average-variance]\label{ass:nondegenerate-average-variance}
\(V_B\to V>0\).
\par\smallskip\noindent\emph{Standard} (nondegenerate limiting design variance, \cite{LiDing2017FPCLT}).
\end{assumption}

\begin{assumption}[ass:bound-lln]\label{ass:bound-lln}
The row-average design variance of the realized bound tables satisfies \(B^{-2}\sum_{j=1}^B\operatorname{Var}(G_{Bj})\longrightarrow0\).
\par\smallskip\noindent\emph{Standard} (triangular-array variance negligibility, \cite{Billingsley1995ProbabilityMeasure}).
\end{assumption}

\begin{assumption}[ass:block-conservativeness]\label{ass:block-conservativeness}
\(b_{Bj}\ge\sigma_{Bj}^2\) for every \(B\) and \(j\).
\par\smallskip\noindent\emph{Standard} (conservative expected randomization-variance bound, \cite{HarshawMiddletonSavje2026}).
\end{assumption}

\section{Definitions}

\begin{definition}[def:observable-complex]\label{def:observable-complex}
\par\smallskip\noindent\emph{Defined object.} \texttt{observable-\allowbreak{}subset complex}
\par\smallskip\noindent\emph{Construction.} \(\mathcal J:=\{S\subseteq I:\exists z\in Z,\ S\subseteq O_z\}\).
\end{definition}

\begin{definition}[def:boolean-mobius-spaces]\label{def:boolean-mobius-spaces}
\par\smallskip\noindent\emph{Defined object.} \texttt{Boolean observable spaces}
\par\smallskip\noindent\emph{Construction.} \(\mathcal E_{\mathcal J}:=\operatorname{span}\{m_S:S\in\mathcal J\}\) and \(\mathcal E_{\mathcal J,r}:=\operatorname{span}\{m_S:S\in\mathcal J,\ |S|\le r\}\), where \(m_S(\theta)=\prod_{i\in S}\theta_i\).
\end{definition}

\begin{definition}[def:score-variance]\label{def:score-variance}
\par\smallskip\noindent\emph{Defined object.} \texttt{linear target and true randomization variance}
\par\smallskip\noindent\emph{Construction.} \(X_z(\theta):=v_z^\top\theta\), \(c:=\sum_zp_zv_z\), \(\tau(\theta):=c^\top\theta\), \(A:=\sum_zp_zv_zv_z^\top-cc^\top\), and \(f(\theta):=\theta^\top A\theta\).
\par\smallskip\noindent\emph{Inputs.} \texttt{ass:\allowbreak{}observable-\allowbreak{}score}.
\end{definition}

\begin{definition}[def:expected-rule]\label{def:expected-rule}
\par\smallskip\noindent\emph{Defined object.} \texttt{expected observed-\allowbreak{}data rule}
\par\smallskip\noindent\emph{Construction.} For tables \(g_z:\{0,1\}^{O_z}\to\mathbb R\), set \(b_g(\theta):=\sum_zp_zg_z(\theta|_{O_z})\).
\end{definition}

\begin{definition}[def:bound-program]\label{def:bound-program}
\par\smallskip\noindent\emph{Defined object.} \texttt{observable-\allowbreak{}margin bound programs}
\par\smallskip\noindent\emph{Construction.} \(\mathcal C_{\mathcal J}(f):=\{b\in\mathcal E_{\mathcal J}:b(\theta)\ge f(\theta)\ \forall\theta\in\Theta\}\), \(F_{\mathcal J}:=\min_{b\in\mathcal C_{\mathcal J}(f)}\sum_\theta q_\theta b(\theta)\), and \(F_{\mathcal J,r}:=\min_{b\in\mathcal C_{\mathcal J}(f)\cap\mathcal E_{\mathcal J,r}}\sum_\theta q_\theta b(\theta)\).
\end{definition}

\begin{definition}[def:dual-and-admissibility]\label{def:dual-and-admissibility}
\par\smallskip\noindent\emph{Defined object.} \texttt{observable-\allowbreak{}margin dual and Pareto order}
\par\smallskip\noindent\emph{Construction.} Set \(\mathcal D_{\mathcal J}(q):=\{\mu\ge0:\sum_\theta\mu_\theta m_S(\theta)=\sum_\theta q_\theta m_S(\theta)\ \forall S\in\mathcal J\}\). Call \(b\in\mathcal C_{\mathcal J}(f)\) Pareto-admissible when no \(b'\in\mathcal C_{\mathcal J}(f)\) satisfies \(b'\le b\) pointwise with a strict inequality somewhere.
\end{definition}

\begin{definition}[def:aronow-samii-binary-correction]\label{def:aronow-samii-binary-correction}
\par\smallskip\noindent\emph{Defined object.} \texttt{binary Aronow-\allowbreak{}-\allowbreak{}Samii zero-\allowbreak{}pairwise-\allowbreak{}inclusion correction}
\par\smallskip\noindent\emph{Construction.} In the Horvitz--Thompson sampling-total specialization with \(\pi_{\{i\}}>0\) and \(v_{z,i}=\mathbb 1\{i\in O_z\}/\pi_{\{i\}}\), define the simplified \(a_{ik}=b_{ik}=2\) correction of Section 3.2 of \cite{AronowSamii2013} by \(g_z^{\mathrm{AS}}(\theta|_{O_z}):=\sum_{i,k\in I:\pi_{\{i,k\}}>0}\mathbb 1\{\{i,k\}\subseteq O_z\}[(\pi_{\{i,k\}}-\pi_{\{i\}}\pi_{\{k\}})\theta_i\theta_k]/[\pi_{\{i,k\}}\pi_{\{i\}}\pi_{\{k\}}]+\sum_{i\ne k:\pi_{\{i,k\}}=0}\{\mathbb 1\{i\in O_z\}\theta_i/[2\pi_{\{i\}}]+\mathbb 1\{k\in O_z\}\theta_k/[2\pi_{\{k\}}]\}\), and set \(b_{\mathrm{AS}}(\theta):=\sum_zp_zg_z^{\mathrm{AS}}(\theta|_{O_z})\).
\end{definition}

\begin{definition}[def:witness-experiment]\label{def:witness-experiment}
\par\smallskip\noindent\emph{Defined object.} \texttt{six-\allowbreak{}coordinate three-\allowbreak{}atom witness and certificates}
\par\smallskip\noindent\emph{Construction.} Let \(\mathfrak E_\star\) have \(K=6\), \(p_z=1/3\), \(O_1=(1,3,5)\), \(O_2=(3,4)\), \(O_3=(0,1,2,5)\), \(v_1=(0,-2,0,-1,0,3)\), \(v_2=(0,0,0,1,-1,0)\), and \(v_3=(-2,-1,2,0,0,1)\). For \(x_i=2\theta_i-1\), put \(P(x)=3x_0x_1-35x_0x_2+3x_0x_5-3x_1x_2+24x_1x_3-45x_1x_5-3x_2x_5-12x_3x_4-36x_3x_5\), \(b^\star(\theta)=[107+P(x)-3x_0x_1x_2x_5]/72\), and \(r^\star(\theta)=[110+P(x)]/72\). Define \(h_1^\star=(0,5,1,0,0,1,5,0)\), \(h_2^\star=(0,1,1,0)\), and \(h_3^\star=(0,2,11/3,14/3,2,0,14/3,8/3,8/3,14/3,0,2,14/3,11/3,2,0)\), in lexicographic order on each \(O_z\). Let \(\mu^\star\) put mass \(1/16\) on \(000000,000111,001011,001110,010011,010100,011100,011111,100000,100111,101011,101100,110000,110101,111000,111011\), and let \(\nu^\star\) put mass \(1/8\) on \(000011,001010,010000,011111,100100,101111,110001,111100\).
\end{definition}

\begin{definition}[def:product-experiment]\label{def:product-experiment}
\par\smallskip\noindent\emph{Defined object.} \texttt{independent product experiment}
\par\smallskip\noindent\emph{Construction.} For disjoint component schedule spaces and independent assignment laws, aggregate local scores as \(\sum_{j=1}^B a_jX_j\), use product objective weights, and retain every joint observed-data table generated by the product observation sets.
\end{definition}

\begin{definition}[def:wald-interval]\label{def:wald-interval}
\par\smallskip\noindent\emph{Defined object.} \texttt{conservative studentized interval}
\par\smallskip\noindent\emph{Construction.} \(I_B(\alpha):=[\widehat\tau_B-z_\alpha\sqrt{\max(\widehat W_B,0)/B},\ \widehat\tau_B+z_\alpha\sqrt{\max(\widehat W_B,0)/B}]\).
\end{definition}

\begin{definition}[def:school-handle]\label{def:school-handle}
\par\smallskip\noindent\emph{Defined object.} \texttt{school-\allowbreak{}network exact-\allowbreak{}decomposition handle}
\par\smallskip\noindent\emph{Construction.} \texttt{Starting from the published two-\allowbreak{}stage school assignment law,\allowbreak{} exposure mapping,\allowbreak{} and direct-\allowbreak{}effect score,\allowbreak{} quotient assignments and observable subsets by exact design automorphisms,\allowbreak{} introduce separator-\allowbreak{}local Boolean Mobius states,\allowbreak{} and require a rational lifting map that preserves every global domination inequality and observable margin;\allowbreak{} the solver determines whether a lossless bounded-\allowbreak{}width decomposition and a useful block exist.}
\end{definition}

\section{Main results}

\begin{theorem}[thm:observable-mobius]\label{thm:observable-mobius}
For any finite randomization with arbitrary real probabilities on its support \(Z\), an expected rule \(b_g:\Theta\to\mathbb R\) is generated by assignment-specific observed-data tables if and only if \(b_g\in\mathcal E_{\mathcal J}\). Equivalently, its unique coefficients \(\beta_S=\sum_{T\subseteq S}(-1)^{|S|-|T|}b_g(\mathbf 1_T)\) vanish for every \(S\notin\mathcal J\). Moreover, every \(b=\sum_{S\in\mathcal J}\beta_Sm_S\) has the assignment-level unbiased implementation \(g_z(\theta|_{O_z})=\sum_{S\subseteq O_z}(\beta_S/\pi_S)m_S(\theta)\), and \(b_g\) is universally conservative exactly when \(b_g(\theta)\ge f(\theta)\) for every \(\theta\in\Theta\).
\end{theorem}
\begin{proof}
Because \(Z\) is the support of the randomization law, \(p_z>0\) for every \(z\in Z\), and \(\sum_{z\in Z}p_z=1\).  The schedule is fixed, and the assignment is the only random object.  In particular, \(\text{def:score-variance}\) keeps the causal target separate from the variance-bound functional:
\[
 \mathbb E_pX_z(\theta)=\sum_zp_zv_z^\top\theta
 =c^\top\theta=\tau(\theta),\qquad
 \operatorname{Var}_p\{X_z(\theta)\}=f(\theta).
\]
First consider an arbitrary function \(u:\Theta\to\mathbb R\).  For \(S\subseteq I\), define
\[
 \gamma_S:=\sum_{T\subseteq S}(-1)^{|S|-|T|}u(\mathbf 1_T).
\]
If \(R\subseteq I\), then finite rearrangement gives
\[
 \sum_{S\subseteq R}\gamma_S
 =\sum_{T\subseteq R}u(\mathbf 1_T)
   \sum_{S:T\subseteq S\subseteq R}(-1)^{|S|-|T|}
 =u(\mathbf 1_R),
\]
because the inner sum equals \((1-1)^{|R\setminus T|}\), interpreted as one when \(T=R\).  Every \(\theta\in\Theta\) equals \(\mathbf 1_R\) for \(R=\{i:\theta_i=1\}\), so
\[
 u(\theta)=\sum_{S\subseteq I}\gamma_Sm_S(\theta).
\]
This expansion is unique: evaluating a zero expansion successively at \(\mathbf 1_R\), beginning with \(R=\varnothing\), forces every coefficient to be zero.  Thus \(\gamma_S\) is the claimed \(\beta_S\).

For necessity, expand each finite table uniquely on its own cube:
\[
 g_z(y)=\sum_{S\subseteq O_z}\gamma_{z,S}m_S(y).
\]
Using \(\operatorname{supp}(v_z)\subseteq O_z\) is unnecessary for this algebra; it is used by \(\text{def:score-variance}\) to ensure that the score itself is observed.  By \(\text{def:expected-rule}\),
\[
 b_g(\theta)=\sum_{z\in Z}p_z\sum_{S\subseteq O_z}
       \gamma_{z,S}m_S(\theta).
\]
Every support occurring on the right belongs to \(\mathcal J\) by \(\text{def:observable-complex}\).  Uniqueness of the Boolean expansion therefore implies \(\beta_S=0\) for \(S\notin\mathcal J\), hence \(b_g\in\mathcal E_{\mathcal J}\) by \(\text{def:boolean-mobius-spaces}\).

Conversely, let \(b=\sum_{S\in\mathcal J}\beta_Sm_S\).  For \(S\in\mathcal J\), define
\[
 \pi_S:=\sum_{z:S\subseteq O_z}p_z.
\]
The defining witness in \(\mathcal J\) and positivity on the support give \(\pi_S>0\); in particular \(\pi_\varnothing=1\).  Hence the displayed assignment rule is well defined.  Its expectation is
\[
 \begin{aligned}
 b_g(\theta)
 &=\sum_zp_z\sum_{S\subseteq O_z}\frac{\beta_S}{\pi_S}m_S(\theta)\\
 &=\sum_{S\in\mathcal J}\frac{\beta_S}{\pi_S}m_S(\theta)
     \sum_{z:S\subseteq O_z}p_z
 =\sum_{S\in\mathcal J}\beta_Sm_S(\theta)=b(\theta).
 \end{aligned}
\]
All sums are finite, so the rearrangement is exact.  Finally, for each fixed schedule, \(\text{def:score-variance}\) gives
\[
 \operatorname{Var}_p\{X_z(\theta)\}=f(\theta),
 \quad\quad \mathbb E_p\{g_z(\theta|_{O_z})\}=b_g(\theta).
\]
Thus expectation conservativeness simultaneously for every fixed binary schedule is equivalent, in both directions, to the finite family of pointwise inequalities \(b_g(\theta)\ge f(\theta)\) for all \(\theta\in\Theta\).
\end{proof}
\par\smallskip\noindent \textit{Justification.} This identifies the complete design-estimable expected-bound space rather than a chosen algebraic subclass. \textit{Gap.} Proposition 1 of \cite{HarshawMiddletonSavje2026} is the degree-two pairwise case; it does not characterize arbitrary observed-data truth tables on binary schedules. \textit{Consumer.} It converts any finite binary exposure experiment into a complete finite set of estimability and conservativeness constraints.

\begin{theorem}[thm:dual-complete-class]\label{thm:dual-complete-class}
For every full-support \(q\), the primal optimum is attained and obeys \(F_{\mathcal J}=\max_{\mu\in\mathcal D_{\mathcal J}(q)}\sum_\theta\mu_\theta f(\theta)\). A feasible \(b\) and feasible \(\mu\) are jointly optimal if and only if \(\mu_\theta>0\) implies \(b(\theta)=f(\theta)\). Every minimizer is Pareto-admissible, and every Pareto-admissible member of \(\mathcal C_{\mathcal J}(f)\) minimizes the displayed objective for some full-support \(q\). These conclusions hold for arbitrary real assignment probabilities on the finite support \(Z\).
\end{theorem}
\begin{proof}
Let \(D=|\mathcal J|\), index the columns by \(S\in\mathcal J\), and define the \(|\Theta|\times D\) matrix and objective vector
\[
 M_{\theta S}:=m_S(\theta),\quad\quad a:=M^\top q.
\]
By the unique expansion proved in \(\text{thm:observable-mobius}\), a coefficient vector \(\beta\in\mathbb R^D\) represents \(b=M\beta\), and \(\text{def:bound-program}\) becomes
\[
 \min_{\beta\in\mathbb R^D}a^\top\beta
 \quad\text{subject to}\quad M\beta\ge f. \tag{1}
\]
The program is feasible because the observable constant \(C:=\max_{\theta\in\Theta}f(\theta)\) dominates \(f\).  Moreover \(f(\theta)\ge0\), since \(f\) is the design variance in \(\text{def:score-variance}\).  If a feasible \(b\) has \(\sum_\theta q_\theta b(\theta)\le L\), full support of \(q\) gives
\[
 0\le b(\theta)\le \frac{L}{q_\theta}\quad\quad(\theta\in\Theta). \tag{2}
\]
The corresponding sublevel set is therefore closed and bounded in the finite-dimensional subspace \(\mathcal E_{\mathcal J}\), hence compact.  Taking \(L=C\) proves primal attainment.

We next derive, rather than assume, the dual certificate.  Let \(\beta^\star\) solve (1), let
\[
 I_0:=\{\theta:M_{\theta\cdot}\beta^\star=f(\theta)\},
 \quad\quad
 \mathcal K:=\operatorname{cone}\{M_{\theta\cdot}^{\top}:\theta\in I_0\}.
\]
We claim that \(a\in\mathcal K\).  Otherwise, choose a closest point \(u\in\mathcal K\) to \(a\); this minimizer exists because a finitely generated cone is closed and the squared distance is coercive after restriction to a sufficiently large ball.  Put \(d=a-u\).  Comparing \(u\) with \(0\), \(2u\), and an arbitrary \(v\in\mathcal K\) in the closest-point inequality gives
\[
 d^\top u=0,\quad\quad d^\top v\le0\quad(v\in\mathcal K).
\]
Consequently \(M_{\theta\cdot}(-d)\ge0\) for every \(\theta\in I_0\), while
\[
 a^\top(-d)=-(u+d)^\top d=-\|d\|_2^2<0. \tag{3}
\]
Every inactive constraint has strictly positive slack.  Finiteness therefore permits an \(\varepsilon>0\) small enough that \(\beta^\star-\varepsilon d\) remains feasible.  Equation (3) would strictly lower the objective, a contradiction.  Thus there are coefficients \(\mu_\theta\ge0\), supported on \(I_0\), such that
\[
 M^\top\mu=a=M^\top q. \tag{4}
\]
The column \(S=\varnothing\) in (4) gives \(\sum_\theta\mu_\theta=\sum_\theta q_\theta=1\), so \(\mu\in\mathcal D_{\mathcal J}(q)\).  Because \(\mu\) is supported on active constraints,
\[
 \sum_\theta\mu_\theta f(\theta)
 =\mu^\top M\beta^\star
 =a^\top\beta^\star=F_{\mathcal J}. \tag{5}
\]
Conversely, for every primal-feasible \(\beta\) and every \(\mu\in\mathcal D_{\mathcal J}(q)\),
\[
 a^\top\beta=\mu^\top M\beta\ge\mu^\top f. \tag{6}
\]
Equations (5)--(6) prove strong duality and dual attainment directly from the finite constraints.  They also show that the primal--dual gap is
\[
 \sum_{\theta\in\Theta}\mu_\theta\{b(\theta)-f(\theta)\}. \tag{7}
\]
All summands in (7) are nonnegative.  Therefore a feasible pair is jointly optimal if and only if \(\mu_\theta>0\) implies \(b(\theta)=f(\theta)\).

The moment equations have the claimed margin interpretation.  For a cell \(y\in\{0,1\}^{O_z}\), its indicator is
\[
 \prod_{i\in O_z:y_i=1}\theta_i
 \prod_{i\in O_z:y_i=0}(1-\theta_i),
\]
which expands into monomials \(m_S\) with \(S\subseteq O_z\), all in \(\mathcal J\).  Hence (4) implies equality of every \(O_z\)-cell probability.  Conversely, equality of each such margin implies equality of the expectation of every \(m_S\) with \(S\subseteq O_z\); the definition of \(\mathcal J\) then gives all equations in (4).

It remains to prove the complete-class assertion.  If \(b\) minimizes an objective with \(q_\theta>0\) and \(b'\le b\) with strict inequality at some schedule, then \(\sum_\theta q_\theta b'(\theta)<\sum_\theta q_\theta b(\theta)\).  Thus every minimizer is Pareto-admissible.

Conversely, fix a Pareto-admissible \(b\) and put
\[
 I_0:=\{\theta:b(\theta)=f(\theta)\},\quad\quad
 \mathcal T:=\{d\in\mathcal E_{\mathcal J}:d(\theta)\ge0\text{ for }\theta\in I_0\},
\]
and
\[
 \mathcal N:=\left\{d\in\mathbb R^\Theta:d\le0,
              \ \sum_\theta d(\theta)=-1\right\}.
\]
Every \(d\in\mathcal T\) is a feasible one-sided direction: on active coordinates it does not reduce the bound, and the finitely many inactive positive slacks allow \(b+\varepsilon d\ge f\) for all sufficiently small \(\varepsilon>0\).  Hence Pareto-admissibility implies \(\mathcal T\cap\mathcal N=\varnothing\).  The set \(\mathcal T\) is a closed polyhedral cone and \(\mathcal N\) is compact.  Choose a closest pair \((t_0,n_0)\) and set \(w=t_0-n_0\ne0\).  The closest-point inequalities, together with the conic comparisons \(0,2t_0\in\mathcal T\), yield
\[
 w^\top t\ge0\quad(t\in\mathcal T),\quad\quad
 w^\top n\le w^\top n_0=-\|w\|_2^2<0\quad(n\in\mathcal N). \tag{8}
\]
Since \(-e_\theta\in\mathcal N\), (8) gives \(w_\theta\ge\|w\|_2^2>0\) for every \(\theta\).  Normalize \(q_\theta=w_\theta/\sum_\eta w_\eta\).  This \(q\) has full support.  For any feasible \(b'\), the difference \(b'-b\) belongs to \(\mathcal T\), so (8) gives
\[
 \sum_\theta q_\theta b'(\theta)\ge
 \sum_\theta q_\theta b(\theta).
\]
Thus every Pareto-admissible bound is supported by a full-support objective.  No rationality was used anywhere in the argument; only finiteness, positivity on the assignment support, and full support of \(q\) were used.
\end{proof}
\par\smallskip\noindent \textit{Justification.} The margin dual supplies exact certificates and the supporting-objective converse supplies the full Pareto complete class. \textit{Gap.} Theorem 3 of \cite{HarshawMiddletonSavje2026} characterizes admissible quadratic matrices through positive-definite matrix objectives; the proposed theorem characterizes all observable Boolean functions through full-support schedule weights and pointwise domination. \textit{Consumer.} A designer can enumerate all undominated expected bounds through supporting weights, without imposing rationality on the experiment.

\begin{proposition}[prop:rational-certificates]\label{prop:rational-certificates}
If every \(p_z\), every coordinate of \(v_z\), and every \(q_\theta\) is rational, then the observable-margin primal and dual programs admit rational optimizers. Finite rational feasibility, equality of the two rational objectives, and complementary slackness form an exact primal--dual certificate; the same implementation formula yields rational assignment-level tables.
\end{proposition}
\begin{proof}
Assume the rationality conditions in the proposition.  By \(\text{def:score-variance}\), each entry of
\[
 c=\sum_zp_zv_z,\quad\quad
 A=\sum_zp_zv_zv_z^\top-cc^\top
\]
is rational, and hence every value \(f(\theta)=\theta^\top A\theta\) is rational.  In the notation of the proof of \(\text{thm:dual-complete-class}\), the matrix \(M\) is a zero--one matrix and \(a=M^\top q\) is rational.  Thus both finite programs are rational polyhedra with rational objectives.

For completeness, rational optimizers follow by elementary active-constraint algebra.  A linear functional on a nonempty compact polytope attains its optimum at an extreme point: if an optimizer is not extreme, pass to an endpoint of a nontrivial segment in the optimal face and iterate within the strictly decreasing face dimension.  At an extreme point, a maximal linearly independent collection of active equalities uniquely determines the point.  Those equalities have rational coefficients and right-hand sides, so Gaussian elimination gives rational coordinates.  The dual feasible set is compact because its empty-monomial equation is \(\sum_\theta\mu_\theta=1\).  It therefore has a rational extreme-point optimizer \(\mu^\star\), and \((\mu^\star)^\top f\) is rational.  By duality this is the primal value.  The primal optimal face is compact by inequality (2) in the proof of \(\text{thm:dual-complete-class}\); applying the same active-equality argument gives a rational coefficient optimizer \(\beta^\star\).

Finite primal feasibility is the list
\[
 (M\beta^\star)_\theta-f(\theta)\ge0
 \quad(\theta\in\Theta),
\]
and finite dual feasibility is
\[
 \mu^\star\ge0,\quad\quad M^\top\mu^\star=M^\top q.
\]
When their rational objectives agree, weak duality proves optimality; equivalently, the rational nonnegative terms in
\[
 \sum_\theta\mu^\star_\theta
       \{(M\beta^\star)_\theta-f(\theta)\}=0
\]
give complementary slackness.  These are exact arithmetic checks with no limiting tolerance.  Finally, Boolean inversion expresses every \(\beta_S^\star\) rationally, and each
\[
 \pi_S=\sum_{z:S\subseteq O_z}p_z
\]
is positive rational for \(S\in\mathcal J\).  The implementation in \(\text{thm:observable-mobius}\) therefore gives rational assignment tables \(g_z\).
\end{proof}
\par\smallskip\noindent \textit{Justification.} Rationality is needed for machine-checkable exact certificates, not for the observable-function or Pareto characterization. \textit{Gap.} The exact finite certificate is a computational corollary of the all-observable program, whereas \cite{HarshawMiddletonSavje2026} certifies a matrix-quadratic program. \textit{Consumer.} Auditors can verify a reported finite optimum using rational arithmetic with no numerical tolerance.

\begin{proposition}[prop:variance-estimability]\label{prop:variance-estimability}
The true variance \(f\) is design-estimable if and only if \(f\in\mathcal E_{\mathcal J}\). In the linear-score model this holds if and only if \(c_ic_k=0\) for every distinct pair \(i,k\) that is never jointly contained in any \(O_z\).
\end{proposition}
\begin{proof}
The first equivalence is exactly \(\text{thm:observable-mobius}\) applied to the function \(f\): an unbiased observed-data variance rule exists if and only if \(f\in\mathcal E_{\mathcal J}\).  On the binary cube, \(\theta_i^2=\theta_i\), and symmetry of \(A\) gives the unique multilinear expansion
\[
 f(\theta)=\sum_{i\in I}A_{ii}\theta_i
       +2\sum_{i<k}A_{ik}\theta_i\theta_k. \tag{9}
\]
If a coordinate \(i\) is never observed, \(\text{ass:observable-score}\) forces \(v_{z,i}=0\) for every \(z\), whence \(c_i=0\), \(A_{ii}=0\), and \(A_{ik}=0\) for every \(k\).  Thus no unobservable singleton occurs in (9).

For distinct \(i,k\) that are never jointly contained in an \(O_z\), observability of the score gives \(v_{z,i}v_{z,k}=0\) for every \(z\).  The definition of \(A\) then yields
\[
 A_{ik}=\sum_zp_zv_{z,i}v_{z,k}-c_ic_k=-c_ic_k. \tag{10}
\]
Every pair that is jointly observed belongs to \(\mathcal J\), while a never-jointly-observed pair does not.  By uniqueness of (9), \(f\in\mathcal E_{\mathcal J}\) holds if and only if \(A_{ik}=0\) for every such distinct pair.  Equation (10) makes this equivalent to \(c_ic_k=0\) for every never-jointly-observed distinct pair, as claimed.  The argument also covers \(K=0\) and \(K=1\), when the pair condition is vacuous.
\end{proof}
\par\smallskip\noindent \textit{Justification.} This is the exact no-slack frontier and a sanity reduction to the familiar pairwise obstruction for a quadratic target. \textit{Gap.} \cite{Middleton2021Unifying} and \cite{HarshawMiddletonSavje2026} give the matrix compatibility criterion over real schedules; the proposition derives its binary Boolean reduction inside the complete class. \textit{Consumer.} It tells experimenters when the LP is unnecessary because an unbiased variance estimator already exists.

\begin{proposition}[prop:hms-class-inclusion]\label{prop:hms-class-inclusion}
Let \(H\) be symmetric and satisfy \(H_{ii}=0\) whenever no \(z\in Z\) has \(i\in O_z\), as well as the published Harshaw--Middleton--Savje pairwise design-compatibility criterion \(H_{ik}=0\) whenever \(i\ne k\) and no \(z\in Z\) has \(\{i,k\}\subseteq O_z\). Then \(h_H\in\mathcal E_{\mathcal J,2}\subseteq\mathcal E_{\mathcal J,3}\). If also \(H-A\succeq0\), then \(h_H\in\mathcal C_{\mathcal J}(f)\cap\mathcal E_{\mathcal J,2}\). The inclusion of published Harshaw--Middleton--Savje conservative quadratic candidates satisfying these observability conditions in the Boolean quadratic cone is strict: the constant function with value \(1+\max_{\theta\in\Theta}f(\theta)\) belongs to \(\mathcal C_{\mathcal J}(f)\cap\mathcal E_{\mathcal J,2}\) but cannot equal any \(h_H\), since every homogeneous quadratic satisfies \(h_H(0)=0\). Consequently the inclusion in the Boolean quadratic/cubic comparison cone is also strict.
\end{proposition}
\begin{proof}
Binary idempotence and symmetry give, for every \(\theta\in\Theta\),
\[
h_H(\theta)=\theta^\top H\theta
=\sum_{i\in I}H_{ii}\theta_i+2\sum_{i<k}H_{ik}\theta_i\theta_k. \tag{1}
\]
If \(H_{ii}\ne0\), the diagonal observability condition supplies an atom \(z\in Z\) with \(i\in O_z\), whence \(\{i\}\in\mathcal J\) by \(\text{def:observable-complex}\).  If \(i\ne k\) and \(H_{ik}\ne0\), pairwise design compatibility supplies an atom \(z\in Z\) with \(\{i,k\}\subseteq O_z\), whence \(\{i,k\}\in\mathcal J\).  Thus every nonzero monomial in (1) is observable and has degree at most two.  By \(\text{def:boolean-mobius-spaces}\),
\[
h_H\in\mathcal E_{\mathcal J,2}\subseteq\mathcal E_{\mathcal J,3}. \tag{2}
\]
If also \(H-A\succeq0\), then the definition of positive semidefiniteness and \(\text{def:score-variance}\) imply
\[
h_H(\theta)-f(\theta)=\theta^\top(H-A)\theta\ge0
\qquad(\theta\in\Theta). \tag{3}
\]
Equations (2)--(3) and \(\text{def:bound-program}\) give
\(h_H\in\mathcal C_{\mathcal J}(f)\cap\mathcal E_{\mathcal J,2}\).

It remains to prove that this inclusion is strict.  Since \(\Theta\) is finite, the number
\[
C:=1+\max_{\theta\in\Theta}f(\theta)
\]
is finite.  Since \(f\) is a variance by \(\text{def:score-variance}\), \(f(\theta)\ge0\), so \(C\ge1\).  The constant function \(b_C(\theta):=C\) is a multiple of \(m_\varnothing\), hence belongs to \(\mathcal E_{\mathcal J,2}\), and \(b_C\ge f\), hence belongs to \(\mathcal C_{\mathcal J}(f)\).  On the other hand, every homogeneous quadratic satisfies \(h_H(0)=0\), while \(b_C(0)=C>0\).  Therefore \(b_C\ne h_H\) for every \(H\).  This proves strict inclusion in the Boolean quadratic cone, and therefore also strict inclusion in the Boolean degree-at-most-three cone.
\end{proof}
\par\smallskip\noindent \textit{Justification.} This embeds every published design-compatible positive-semidefinite quadratic candidate into the Boolean quadratic/cubic comparison cone while keeping the strict class distinction explicit. \textit{Gap.} Proposition 1 and Definitions 1--2 of \cite{HarshawMiddletonSavje2026} give the pairwise matrix criterion and positive-semidefinite conservativeness; they do not state the strict binary-cube inclusion or solve the present witness program. \textit{Consumer.} The inclusion licenses the \(55/36\) enlarged-class lower bound, while the separate exact semidefinite certificate sharpens the direct published-class comparison to \(7/6+\sqrt{17}/9\).

\begin{theorem}[thm:complement-degree-collapse]\label{thm:complement-degree-collapse}
Under the proposal-specific conditions \(A\mathbf 1_K=0\) and \(q_\theta=q_{\mathbf 1_K-\theta}\), \(F_{\mathcal J,2d+1}=F_{\mathcal J,2d}\) for every integer \(d\ge0\). In particular, cubic terms cannot improve the quadratic optimum, and if \(|O_z|\le3\) for every \(z\), the quadratic optimum equals the all-observable optimum. When the objective weights are not complement invariant, symmetrization need not preserve the objective and this theorem does not rule out a degree-three improvement.
\end{theorem}
\begin{proof}
The matrix \(A\) in \(\text{def:score-variance}\) is symmetric.  Under \(A\mathbf 1_K=0\), for every \(\theta\in\Theta\),
\[
 \begin{aligned}
 f(\mathbf 1_K-\theta)
 &=(\mathbf 1_K-\theta)^\top A(\mathbf 1_K-\theta)\\
 &=\mathbf 1_K^\top A\mathbf 1_K
   -2\theta^\top A\mathbf 1_K+\theta^\top A\theta
 =f(\theta).
 \end{aligned}\tag{13}
\]
If \(b\) is feasible, define its complement symmetrization by
\[
 \bar b(\theta):=\frac{b(\theta)+b(\mathbf 1_K-\theta)}2.
\]
Equation (13) shows \(\bar b\ge f\).  Moreover, complementation preserves \(\mathcal E_{\mathcal J}\): for every \(S\in\mathcal J\),
\[
 m_S(\mathbf 1_K-\theta)=\prod_{i\in S}(1-\theta_i)
 =\sum_{T\subseteq S}(-1)^{|T|}m_T(\theta),
\]
and \(T\in\mathcal J\) by downward closure.  Complement invariance of \(q\) gives, after the bijective change of variable \(\eta=\mathbf 1_K-\theta\),
\[
 \sum_\theta q_\theta\bar b(\theta)
 =\frac12\sum_\theta q_\theta b(\theta)
  +\frac12\sum_\eta q_{\mathbf 1_K-\eta}b(\eta)
 =\sum_\theta q_\theta b(\theta). \tag{14}
\]

Now use \(x_i=2\theta_i-1\).  For each \(S\), the products
\(\prod_{i\in S}x_i\) and \(m_S\) are related by an invertible triangular change of basis involving only subsets of \(S\); hence this basis preserves both membership in \(\mathcal J\) and degree.  Complementation sends \(x\) to \(-x\).  Therefore symmetrization deletes every term of odd cardinality and leaves every even term unchanged.  In particular, a feasible \(b\in\mathcal E_{\mathcal J,2d+1}\) produces, by (13)--(14), a feasible \(\bar b\in\mathcal E_{\mathcal J,2d}\) with the same objective.  Thus
\[
 F_{\mathcal J,2d}\le F_{\mathcal J,2d+1}.
\]
The reverse inequality follows from \(\mathcal E_{\mathcal J,2d}\subseteq\mathcal E_{\mathcal J,2d+1}\), proving equality.  If every \(|O_z|\le3\), every \(S\in\mathcal J\) has \(|S|\le3\), so \(\mathcal E_{\mathcal J}=\mathcal E_{\mathcal J,3}\).  Taking \(d=1\) then gives
\[
 F_{\mathcal J}=F_{\mathcal J,3}=F_{\mathcal J,2}.
\]
When \(q\) is not complement invariant, equation (14) need not hold; accordingly the proof supplies no exclusion of a cubic improvement in that regime.
\end{proof}
\par\smallskip\noindent \textit{Justification.} Equal total score loadings make outcome complementation a variance-preserving design operation; paired schedule weights then identify the structural reason that the first possible gain in the witness occurs at degree four. \textit{Gap.} The binary inclusion-exclusion argument in \cite{Lu2019FactorialBinary} is design-specific, while the quadratic framework of \cite{HarshawMiddletonSavje2026} has no degree filtration or complement-parity theorem. \textit{Consumer.} It rules out fruitless odd-degree searches and gives an exact stopping rule for observation sets of size at most three.

\begin{theorem}[thm:quartic-separation]\label{thm:quartic-separation}
For \(\mathfrak E_\star\) with uniform \(q\), \(F_{\mathcal J}=107/72\) and \(F_{\mathcal J,2}=F_{\mathcal J,3}=55/36\), so the exact gap is \(1/24\). By the strict class-inclusion proposition, every quadratic candidate allowed by the published Harshaw--Middleton--Savje framework lies in the Boolean comparison class \(\mathcal C_{\mathcal J}(f)\cap\mathcal E_{\mathcal J,2}\subseteq\mathcal C_{\mathcal J}(f)\cap\mathcal E_{\mathcal J,3}\); hence the same \(55/36\) lower certificate rules out every such published candidate, although it does not assert that \(55/36\) is the optimum of their smaller positive-semidefinite class. The displayed bounds satisfy \(r^\star(\theta)-b^\star(\theta)=[1+x_0x_1x_2x_5]/24\in\{0,1/12\}\), and \(b^\star(\theta)=3^{-1}\sum_{z=1}^3h_z^\star(\theta|_{O_z})\) with every table entry in \([0,5]\). The distribution \(\mu^\star\) certifies \(107/72\), and \(\nu^\star\) certifies \(55/36\), by exact observable-margin equalities and complementary slackness across all sixty-four schedules.
\end{theorem}
\begin{proof}
All calculations in this proof are in \(\mathbb Q\), except for the already-established algebraic value of the Harshaw--Middleton--Savje semidefinite program.  The three score vectors in Definition~\texttt{def:witness-experiment} have supports contained in their respective observation sets.  Substitution into Definition~\texttt{def:score-variance} gives
\[
c=\left(-\frac23,-1,\frac23,0,-\frac13,\frac43\right)
\tag{1}
\]
and
\[
9A=
\begin{pmatrix}
8&0&-8&0&-2&2\\
0&6&0&6&-3&-9\\
-8&0&8&0&2&-2\\
0&6&0&6&-3&-9\\
-2&-3&2&-3&2&4\\
2&-9&-2&-9&4&14
\end{pmatrix}.
\tag{2}
\]
In particular, every row of \(A\) sums to zero, so \(A\mathbf 1_6=0\).  The observation sets in Definition~\texttt{def:witness-experiment} yield an observable complex consisting of \(\varnothing\), the six singletons, the nine pairs
\[
01,\ 02,\ 05,\ 12,\ 13,\ 15,\ 25,\ 34,\ 35,
\tag{3}
\]
the five triples
\[
012,\ 015,\ 025,\ 125,\ 135,
\tag{4}
\]
and the quartet \(0125\).  Here and below a digit string denotes the corresponding subset of \(I\).

We first give the complete Boolean coefficient certificate for the displayed, symmetrized \(b^\star\).  Expanding \(x_i x_k=(2\theta_i-1)(2\theta_k-1)\) and
\[
x_0x_1x_2x_5=
\prod_{i\in\{0,1,2,5\}}(2\theta_i-1)
\tag{5}
\]
in Definition~\texttt{def:witness-experiment} gives the following coefficient on every member of \(\mathcal J\):
\[
\begin{array}{c|rrrrrrr}
S&\varnothing&0&1&2&3&4&5\\ \hline
\beta_S&0&8/9&2/3&11/9&2/3&1/3&7/3
\end{array}
\tag{6}
\]
\[
\begin{array}{c|rrrrrrrrr}
S&01&02&05&12&13&15&25&34&35\\ \hline
\beta_S&0&-19/9&0&-1/3&4/3&-8/3&-1/3&-2/3&-2
\end{array}
\tag{7}
\]
\[
\begin{array}{c|rrrrrr}
S&012&015&025&125&135&0125\\ \hline
\beta_S&1/3&1/3&1/3&1/3&0&-2/3.
\end{array}
\tag{8}
\]
Every coefficient indexed by \(S\notin\mathcal J\) is zero.  Equations (6)--(8), together with this last zero statement, specify all \(2^6=64\) Boolean Möbius coefficients.

The same certificate follows independently from the twenty-eight displayed assignment-table entries.  Boolean inversion on each local cube gives
\[
\begin{aligned}
h_1^\star
 &=5m_5+m_3-6m_{35}-4m_{15}+4m_{13},\\
h_2^\star
 &=m_4+m_3-2m_{34},\\
h_3^\star
 &=2m_5+\frac{11}{3}m_2-m_{25}+2m_1-4m_{15}-m_{12}+m_{125}\\
 &\qquad+\frac83m_0-\frac{19}{3}m_{02}+m_{025}+m_{015}+m_{012}-2m_{0125}.
\end{aligned}
\tag{9}
\]
Consequently, the coefficients of \(3^{-1}(h_1^\star+h_2^\star+h_3^\star)\) are exactly (6)--(8).  For example, the two singleton coefficients whose values distinguish this symmetrized rule are
\[
\beta_{\{5\}}=\frac{5+2}{3}=\frac73,
\qquad
\beta_{\{4\}}=\frac{1}{3}.
\tag{10}
\]
For completeness, the Boolean expansion is unique: evaluation at \(\mathbf 1_T\) gives \(u(\mathbf 1_T)=\sum_{S\subseteq T}\beta_S\), and induction on \(|T|\) recovers each \(\beta_T\) from these values.  Therefore equality of all coefficients in (6)--(9) proves all sixty-four implementation equalities
\[
b^\star(\theta)=\frac13\sum_{z=1}^3
h_z^\star(\theta|_{O_z}),
\qquad \theta\in\Theta.
\tag{11}
\]
Inspection of the displayed rational entries gives
\[
0\le h_z^\star(y)\le5
\quad
\text{for every }z\in\{1,2,3\}
\text{ and }y\in\{0,1\}^{O_z}.
\tag{12}
\]
Since all supports in (6)--(8) are in \(\mathcal J\), the definition of the Boolean observable space gives \(b^\star\in\mathcal E_{\mathcal J}\).  Since \(\beta_{0125}=-2/3\ne0\), uniqueness also gives \(b^\star\notin\mathcal E_{\mathcal J,3}\).  The formula for \(r^\star\) contains only a constant and the nine observable centered-coordinate pairs in (3), so the same expansion gives \(r^\star\in\mathcal E_{\mathcal J,2}\).

We next verify every pointwise domination constraint.  Order \(\Theta\) lexicographically as \(000000,000001,\ldots,111111\).  Substitution of (2) into \(f(\theta)=\theta^\top A\theta\) and exact subtraction from the displayed formula for \(b^\star\) gives
\[
72\{b^\star-f\}=
\begin{array}{rrrrrrrr}
0&56&8&0&0&56&8&0\\
24&88&0&0&24&88&0&0\\
0&8&56&0&0&8&56&0\\
0&40&24&0&0&40&24&0\\
0&24&40&0&0&24&40&0\\
0&56&8&0&0&56&8&0\\
0&0&88&24&0&0&88&24\\
0&8&56&0&0&8&56&0.
\end{array}
\tag{13}
\]
This table is the complete list of sixty-four rational slacks, and every entry is nonnegative.  Hence \(b^\star\ge f\) on \(\Theta\).  If
\[
w(\theta):=x_0x_1x_2x_5\in\{-1,1\},
\tag{14}
\]
then direct subtraction of the two formulas in Definition~\texttt{def:witness-experiment} yields
\[
r^\star(\theta)-b^\star(\theta)
=\frac{110-107+3w(\theta)}{72}
=\frac{1+w(\theta)}{24}
\in\left\{0,\frac1{12}\right\}.
\tag{15}
\]
The displayed difference is a nondegenerate observable contrast, rather than a purely averaged artifact.  Put
\[
 \mathcal R:=\{(0,0,1,u,v,1):u,v\in\{0,1\}\}.
\tag{15a}
\]
Under the uniform schedule weights, \(q(\mathcal R)=4/64=1/16\), and \(w(\theta)=1\) for every \(\theta\in\mathcal R\).  Thus the four schedules in this positive-mass region all have the pointwise improvement
\[
 r^\star(\theta)-b^\star(\theta)=\frac1{12}.
\tag{15b}
\]
This contrast has a direct assignment-table implementation: take \(\delta_1=\delta_2=0\) and, on \(O_3=\{0,1,2,5\}\), set
\[
 \delta_3(y):=
 \frac{1+(2y_0-1)(2y_1-1)(2y_2-1)(2y_5-1)}8
 \in\left\{0,\frac14\right\}.
\tag{15c}
\]
Because every assignment atom has probability \(p_z=1/3\), hence \(\min_zp_z=1/3\), its expected observed-data contrast is exactly
\[
 \sum_{z=1}^3p_z\delta_z(\theta|_{O_z})
 =\frac13\delta_3(\theta|_{O_3})
 =\frac{1+w(\theta)}{24}
 =r^\star(\theta)-b^\star(\theta).
\tag{15d}
\]
It is therefore observable and nonnegative assignmentwise, and it is strictly positive throughout \(\mathcal R\).  The randomization is also nondegenerate there: in the order \((u,v)=(0,0),(0,1),(1,0),(1,1)\), the three assignment scores are respectively
\[
 (3,0,3),\quad(3,-1,3),\quad(2,1,3),\quad(2,0,3),
\]
so their equal-probability conditional variances are respectively
\[
 2,\quad\frac{32}{9},\quad\frac23,\quad\frac{14}{9}.
\tag{15e}
\]
Consequently \(f(\theta)=\operatorname{Var}_p\{X_z(\theta)\}\ge2/3\) on the same region.  This supplies quantitative positive lower bounds simultaneously for assignment support, schedule mass, conditional variance, and the witness's pointwise gap.

Therefore \(r^\star\ge b^\star\ge f\).  Under uniform \(q\), every nonconstant product of centered coordinates has expectation zero.  The constant terms in the two displayed formulas consequently give
\[
\mathbb E_q b^\star=\frac{107}{72},
\qquad
\mathbb E_q r^\star=\frac{110}{72}=\frac{55}{36}.
\tag{16}
\]

For the all-degree lower certificate, each of the eight \(O_1\)-patterns occurs exactly twice among the sixteen schedules supporting \(\mu^\star\), each of the four \(O_2\)-patterns occurs exactly four times, and each of the sixteen \(O_3\)-patterns occurs exactly once.  Since \(\mu^\star\) is uniform on these sixteen distinct schedules, it is a probability measure and has the uniform \(q\)-margin on every \(O_z\).  Every \(S\in\mathcal J\) is contained in an \(O_z\), so all observable moments agree and \(\mu^\star\in\mathcal D_{\mathcal J}(q)\).  In the support order displayed in Definition~\texttt{def:witness-experiment}, exact evaluation with (2) gives
\[
0,\ \frac23,\ \frac{32}{9},\ \frac{14}{9},\ \frac23,\ \frac83,
\ \frac{32}{9},\ \frac89,\ \frac89,\ \frac{14}{9},\ \frac83,\ \frac23,
\ \frac{14}{9},\ \frac{14}{9},\ \frac23,\ \frac23.
\tag{17}
\]
Their sum is \(214/9\), so their \(\mu^\star\)-mean is \(107/72\).  The entries of (13) at all sixteen support schedules are zero, proving complementary slackness.  The dual equality and optimality criterion in Theorem~\texttt{thm:dual-complete-class}, combined with (16), therefore give
\[
F_{\mathcal J}=\frac{107}{72}.
\tag{18}
\]

For the degree-three lower certificate, the eight schedules supporting \(\nu^\star\) contain every \(O_1\)-pattern once and every \(O_2\)-pattern twice.  Their \(O_3\)-projections are precisely the eight odd-parity strings in \(\{0,1\}^4\).  For a proper subset \(S\subsetneq O_3\), fixing its \(|S|\) bits leaves exactly \(2^{3-|S|}\) odd-parity completions.  Hence every proper-subset monomial has
\[
\mathbb E_{\nu^\star}m_S=2^{-|S|}=\mathbb E_qm_S.
\tag{19}
\]
Together with the \(O_1\)- and \(O_2\)-counts, equation (19) proves exact agreement for all twenty-one members of \(\{S\in\mathcal J:|S|\le3\}\).  The corresponding values of \(f\), again from (2), are
\[
\frac83,\ \frac{14}{9},\ \frac23,\ \frac89,
\ \frac{14}{9},\ \frac23,\ \frac{14}{9},\ \frac83.
\tag{20}
\]
Their sum is \(110/9\), so their mean is \(55/36\); direct substitution gives \(r^\star=f\) at all eight support schedules.  If \(b\in\mathcal C_{\mathcal J}(f)\cap\mathcal E_{\mathcal J,3}\), moment agreement and pointwise domination imply
\[
\mathbb E_qb=\mathbb E_{\nu^\star}b
\ge\mathbb E_{\nu^\star}f=\frac{55}{36}.
\tag{21}
\]
Feasibility of \(r^\star\), the inclusion \(\mathcal E_{\mathcal J,2}\subseteq\mathcal E_{\mathcal J,3}\), and (21) yield
\[
\frac{55}{36}\le F_{\mathcal J,3}
\le F_{\mathcal J,2}
\le\frac{55}{36}.
\tag{22}
\]
Thus
\[
F_{\mathcal J,2}=F_{\mathcal J,3}=\frac{55}{36},
\qquad
F_{\mathcal J,2}-F_{\mathcal J}=\frac1{24}.
\tag{23}
\]
The unique missing observable moment is the quartet:
\[
\mathbb E_{\nu^\star}m_{0125}=0,
\qquad
\mathbb E_qm_{0125}=\frac1{16},
\tag{24}
\]
equivalently \(\mathbb E_{\nu^\star}x_0x_1x_2x_5=-1\) whereas the uniform \(q\)-expectation is zero.  Moreover, (2) gives \(A\mathbf 1_6=0\), and uniform \(q\) is complement invariant.  Theorem~\texttt{thm:complement-degree-collapse} therefore independently forces \(F_{\mathcal J,3}=F_{\mathcal J,2}\).  Since the complex (3)--(4) has no observable degree above four and its only degree-four member is \(0125\), the quartet term in (8) is the first even degree at which an improvement over degree two can occur.

Finally, every coordinate is observed with positive probability in this witness.  Proposition~\texttt{prop:hms-class-inclusion} therefore places every Harshaw--Middleton--Savje compatible conservative quadratic candidate in \(\mathcal C_{\mathcal J}(f)\cap\mathcal E_{\mathcal J,2}\).  Equation (22) implies that every such candidate has uniform objective at least \(55/36\), without asserting that this is the optimum of the smaller semidefinite class.  Proposition~\texttt{prop:witness-hms-sdp-optimum} supplies the exact optimum of that smaller class:
\[
\min_H\mathbb E_qh_H
=\frac76+\frac{\sqrt{17}}9.
\tag{25}
\]
It follows algebraically that the direct published-class gap and its excess over the enlarged Boolean quadratic optimum are
\[
\left(\frac76+\frac{\sqrt{17}}9\right)-\frac{107}{72}
=\frac{8\sqrt{17}-23}{72}>0,
\qquad
\left(\frac76+\frac{\sqrt{17}}9\right)-\frac{55}{36}
=\frac{4\sqrt{17}-13}{36}>0.
\tag{26}
\]
The inequalities follow by squaring the positive comparisons \(8\sqrt{17}>23\) and \(4\sqrt{17}>13\).  This completes the exact primal, dual, implementation, degree-threshold, and comparator certificates. \(\square\)
\end{proof}
\par\smallskip\noindent \textit{Justification.} The complete rational certificate expands both the centered-coordinate formula and every assignment table, checks all sixty-four slacks and implementation equalities, and distinguishes the all-observable quartic optimum, the larger Boolean quadratic/cubic optimum, and the smaller published positive-semidefinite quadratic optimum. \textit{Gap.} Neither \cite{HarshawMiddletonSavje2026} nor \cite{Lu2019FactorialBinary} contains this quartic witness, the \(1/24\) Boolean degree gap, or the exact \((8\sqrt{17}-23)/72\) direct quadratic-class gap. \textit{Consumer.} It supplies a reproducible rational truth-table test and an exact algebraic semidefinite benchmark without equating the two quadratic comparison classes.

\begin{theorem}[thm:aronow-samii-specialization]\label{thm:aronow-samii-specialization}
In the binary Horvitz--Thompson sampling-total specialization, the simplified Aronow--Samii correction satisfies \(b_{\mathrm{AS}}\in\mathcal C_{\mathcal J}(f)\cap\mathcal E_{\mathcal J,2}\). Thus its weak conservativeness remains valid when some pairwise observation probabilities vanish, and its expected correction is one feasible point of the Boolean observable cone. Across the present model class this embedding is strict: in \(\mathfrak E_\star\), \(b^\star\in\mathcal C_{\mathcal J}(f)\setminus\mathcal E_{\mathcal J,3}\) and improves the degree-at-most-three optimum by \(1/24\).
\end{theorem}
\begin{proof}
Write \(R_i(z)=\mathbb 1\{i\in O_z\}\), \(\pi_i=\mathbb E_pR_i\), and \(\pi_{ik}=\mathbb E_p(R_iR_k)\).  The specialization assumes \(\pi_i>0\) and has
\[
 X_z(\theta)=\sum_i\frac{R_i(z)\theta_i}{\pi_i},
 \quad\quad \mathbb E_pX_z(\theta)=\sum_i\theta_i.
\]
Expanding the variance and using that the schedule is fixed gives
\[
 f(\theta)=\sum_{i,k}
 \frac{\pi_{ik}-\pi_i\pi_k}{\pi_i\pi_k}\theta_i\theta_k. \tag{25}
\]
For \(\pi_{ik}>0\), the expectation of the corresponding first term in \(g_z^{\mathrm{AS}}\) is exactly the summand in (25), because \(\mathbb E_pR_iR_k=\pi_{ik}\).  If \(i\ne k\) and \(\pi_{ik}=0\), the summand in (25) is \(-\theta_i\theta_k\), whereas the expected replacement in \(g_z^{\mathrm{AS}}\) is
\[
 \mathbb E_p\left\{
 \frac{R_i\theta_i}{2\pi_i}+\frac{R_k\theta_k}{2\pi_k}
 \right\}=\frac{\theta_i+\theta_k}{2}. \tag{26}
\]
There are no zero-probability diagonal pairs because \(\pi_{ii}=\pi_i>0\).  Subtracting (25) from the expectation of the displayed Aronow--Samii rule yields
\[
 b_{\mathrm{AS}}(\theta)-f(\theta)
 =\sum_{\substack{i\ne k\\\pi_{ik}=0}}
 \left\{\frac{\theta_i+\theta_k}{2}+\theta_i\theta_k\right\}
 =\frac12\sum_{\substack{i\ne k\\\pi_{ik}=0}}
       (\theta_i+\theta_k)^2\ge0, \tag{27}
\]
where the second equality uses \(\theta_i^2=\theta_i\).  Thus weak conservativeness holds for every fixed binary schedule, including when pairwise observation probabilities vanish.

The expectation of the positive-\(\pi_{ik}\) portion is a linear combination of \(m_{\{i,k\}}\) on jointly observed pairs, with diagonal terms reduced to singletons; (26) contributes only observable singletons.  Since \(p_z>0\) on \(Z\), \(\pi_{ik}>0\) is equivalent to the existence of an assignment observing \(\{i,k\}\).  Therefore
\[
 b_{\mathrm{AS}}\in\mathcal E_{\mathcal J,2},
 \quad\quad b_{\mathrm{AS}}\ge f,
\]
and \(\text{def:bound-program}\) gives
\(b_{\mathrm{AS}}\in\mathcal C_{\mathcal J}(f)\cap\mathcal E_{\mathcal J,2}\).

For strictness across the full finite-binary model class, \(\text{thm:quartic-separation}\) proves in \(\mathfrak E_\star\) that \(b^\star\in\mathcal C_{\mathcal J}(f)\), that its unique quartet coefficient is \(-2/3\), and hence that \(b^\star\notin\mathcal E_{\mathcal J,3}\).  The same theorem gives the exact objective improvement
\[
 F_{\mathcal J,3}-F_{\mathcal J}=\frac{55}{36}-\frac{107}{72}=\frac1{24}.
\]
Thus the published correction remains a feasible quadratic baseline, while the all-observable Boolean cone is strictly larger.
\end{proof}
\par\smallskip\noindent \textit{Justification.} This locates the canonical zero-pairwise-inclusion correction inside the proposed cone and states the strict class enlargement at theorem level. \textit{Gap.} Proposition 2 and Corollary 2 in Section 3 of \cite{AronowSamii2013} prove weak conservativeness of a Horvitz--Thompson correction using Young's inequality; they do not characterize every Boolean observed-data expectation, the Pareto complete class, or a strict quartic-versus-degree-at-most-three frontier. \textit{Consumer.} Survey samplers and exposure-design analysts can retain the published correction as a certified feasible baseline while using the observable-margin program to test and certify higher-order improvements.

\begin{theorem}[thm:tensorization]\label{thm:tensorization}
For the independent product experiment, the full and degree-restricted optima satisfy \(F_{\mathrm{prod}}=\sum_{j=1}^Ba_j^2F_j\) and \(F_{\mathrm{prod},r}=\sum_{j=1}^Ba_j^2F_{j,r}\), even when the comparison class contains joint observed-data interactions across blocks. Primal rules add and dual certificates tensor. For \(B\) identical copies of \(\mathfrak E_\star\) with \(a_j=1/B\), the optima are \(107/(72B)\) and \(55/(36B)\), and the gap is \(1/(24B)\).
\end{theorem}
\begin{proof}
Write the product schedule as \(\theta=(\theta_1,\ldots,\theta_B)\) on disjoint coordinate blocks and the product assignment as \(z=(z_1,\ldots,z_B)\).  All variances below are randomization variances with these schedules fixed.  Definition~\texttt{def:product-experiment} constructs the assignment law as the product of the independent component laws (the \texttt{prodDesign} construction). Hence the block scores \(X_j(\theta_j)\) are independent, every cross-block covariance is zero, and
\[
\begin{aligned}
f_{\mathrm{prod}}(\theta)
&=\operatorname{Var}\!\left(\sum_{j=1}^Ba_jX_j(\theta_j)\right)\\
&=\sum_{j=1}^Ba_j^2\operatorname{Var}\{X_j(\theta_j)\}
=\sum_{j=1}^Ba_j^2f_j(\theta_j),
\end{aligned}
\tag{27}
\]
This is the required product-design variance identity.  Thus the causal target remains the mean of the displayed product score, while (27) is its design variance; neither object is defined to equal a bound.

For each block choose an all-degree primal minimizer \(b_j\), whose existence follows from Theorem~\texttt{thm:dual-complete-class}.  Each lift \(b_j(\theta_j)\) is observable in the product experiment, so
\[
b_{\mathrm{prod}}(\theta):=
\sum_{j=1}^Ba_j^2b_j(\theta_j)
\ge \sum_{j=1}^Ba_j^2f_j(\theta_j)
=f_{\mathrm{prod}}(\theta).
\tag{28}
\]
If \(g_{j,z_j}\) is an assignment-specific observed-data table implementing \(b_j\), then
\[
g_z^{\mathrm{prod}}:=\sum_{j=1}^Ba_j^2g_{j,z_j}
\tag{29}
\]
is a joint observed-data table and its product-design expectation is (28).  Under the product objective weights \(q=\bigotimes_jq_j\),
\[
\mathbb E_qb_{\mathrm{prod}}
=\sum_{j=1}^Ba_j^2\mathbb E_{q_j}b_j
=\sum_{j=1}^Ba_j^2F_j.
\tag{30}
\]
Hence \(F_{\mathrm{prod}}\le\sum_ja_j^2F_j\).  The same construction from degree-at-most-\(r\) local minimizers is still degree at most \(r\), so
\[
F_{\mathrm{prod},r}\le\sum_{j=1}^Ba_j^2F_{j,r}.
\tag{31}
\]

Before proving the restricted reverse inequality, we derive the required restricted duality and attainment rather than assuming them.  For block \(j\), let \(M_{j,r}\) be the finite matrix whose rows are indexed by \(\theta_j\in\Theta_j\), whose columns are indexed by
\[
\mathcal J_{j,r}:=\{S\in\mathcal J_j:|S|\le r\},
\tag{32}
\]
and whose entry is \((M_{j,r})_{\theta_j,S}=m_S(\theta_j)\).  With \(a_{j,r}:=M_{j,r}^{\top}q_j\), the defining local restricted program is exactly
\[
F_{j,r}
=\min_{\beta\in\mathbb R^{\mathcal J_{j,r}}}
\left\{a_{j,r}^{\top}\beta:
M_{j,r}\beta\ge f_j\right\}.
\tag{33}
\]
The empty-set column is constant one, so (33) is feasible by taking a sufficiently large constant.  Pointwise nonnegativity \(f_j\ge0\) and full support of \(q_j\) imply that, on every objective sublevel \(\mathbb E_{q_j}b\le L\),
\[
0\le b(\theta_j)\le \frac{L}{\min_{u\in\Theta_j}q_j(u)}
\quad(\theta_j\in\Theta_j).
\tag{34}
\]
Thus the feasible sublevel is compact in the finite-dimensional value space \(\operatorname{col}(M_{j,r})\), and the restricted primal minimum is attained.

For \(\mu_j\ge0\), the Lagrangian of (33) is
\[
L(\beta,\mu_j)
=\mu_j^{\top}f_j
+\beta^{\top}(a_{j,r}-M_{j,r}^{\top}\mu_j).
\tag{35}
\]
Its infimum over the unrestricted coefficient vector \(\beta\) is finite exactly when
\[
M_{j,r}^{\top}\mu_j=M_{j,r}^{\top}q_j.
\tag{36}
\]
The finite-dimensional primal--dual separation argument used in Theorem~\texttt{thm:dual-complete-class} depends only on the displayed matrix, feasibility, and finite optimal value.  Applying that same argument to \(M_{j,r}\), rather than to the untruncated observable-moment matrix, proves the equality
\[
F_{j,r}
=\max_{\substack{\mu_j\ge0\\
M_{j,r}^{\top}\mu_j=M_{j,r}^{\top}q_j}}
\mu_j^{\top}f_j.
\tag{37}
\]
This is not an invocation of the unrestricted conclusion for a restricted optimizer: equations (32)--(36) reconstruct the restricted primal and its own dual.  The empty-set equation in (36) gives \(\sum_{\theta_j}\mu_j(\theta_j)=1\), so its feasible set is a closed subset of the finite probability simplex.  It is nonempty because \(\mu_j=q_j\), and hence the maximum in (37) is attained.  We may therefore choose a restricted dual optimizer \(\mu_{j,r}\) for every block.  The same argument with all columns gives the all-degree dual optimizers \(\mu_j\) already supplied by Theorem~\texttt{thm:dual-complete-class}.

For the all-degree reverse inequality, put \(\mu:=\bigotimes_{j=1}^B\mu_j\).  Every monomial observable under a product assignment has the unique factorization
\[
m_S(\theta)=\prod_{j=1}^Bm_{S_j}(\theta_j),
\qquad
S_j\subseteq O_{j,z_j}.
\tag{38}
\]
Local observable-moment equality and product factorization imply
\[
\mathbb E_\mu m_S
=\prod_{j=1}^B\mathbb E_{\mu_j}m_{S_j}
=\prod_{j=1}^B\mathbb E_{q_j}m_{S_j}
=\mathbb E_qm_S.
\tag{39}
\]
Thus \(\mu\) is feasible for every product-observable moment, including moments that interact across blocks.  From (27) and local dual optimality,
\[
\mathbb E_\mu f_{\mathrm{prod}}
=\sum_{j=1}^Ba_j^2\mathbb E_{\mu_j}f_j
=\sum_{j=1}^Ba_j^2F_j.
\tag{40}
\]
For every globally feasible bound \(b\), moment equality gives \(\mathbb E_qb=\mathbb E_\mu b\ge\mathbb E_\mu f_{\mathrm{prod}}\).  Equation (40) is therefore a dual lower bound, which together with (30) proves
\[
F_{\mathrm{prod}}=\sum_{j=1}^Ba_j^2F_j.
\tag{41}
\]

For the restricted reverse inequality, put \(\mu_r:=\bigotimes_{j=1}^B\mu_{j,r}\).  If the global monomial in (38) has total degree at most \(r\), then
\[
|S_j|\le\sum_{\ell=1}^B|S_\ell|=|S|\le r
\tag{42}
\]
for every \(j\).  Each factor in (38) is consequently one of the local restricted moments in (36).  Replacing \(\mu_j\) by \(\mu_{j,r}\) in (39) proves that \(\mu_r\) matches \(q\) on every global degree-at-most-\(r\) observable monomial.  Moreover, (27) and (37) give
\[
\mathbb E_{\mu_r}f_{\mathrm{prod}}
=\sum_{j=1}^Ba_j^2\mathbb E_{\mu_{j,r}}f_j
=\sum_{j=1}^Ba_j^2F_{j,r}.
\tag{43}
\]
For every \(b\in\mathcal C_{\mathcal J}(f_{\mathrm{prod}})\cap\mathcal E_{\mathcal J,r}\), restricted moment equality and domination give
\[
\mathbb E_qb=\mathbb E_{\mu_r}b
\ge\mathbb E_{\mu_r}f_{\mathrm{prod}}
=\sum_{j=1}^Ba_j^2F_{j,r}.
\tag{44}
\]
Equations (31) and (44) prove
\[
F_{\mathrm{prod},r}=\sum_{j=1}^Ba_j^2F_{j,r}.
\tag{45}
\]
All formulas remain valid when some \(a_j=0\), since only \(a_j^2\) occurs.

Finally, for \(B\) identical copies of \(\mathfrak E_\star\) with \(a_j=1/B\), Theorem~\texttt{thm:quartic-separation} and equations (41) and (45) give
\[
F_{\mathrm{prod}}
=B\frac1{B^2}\frac{107}{72}
=\frac{107}{72B},
\qquad
F_{\mathrm{prod},2}=F_{\mathrm{prod},3}
=B\frac1{B^2}\frac{55}{36}
=\frac{55}{36B}.
\tag{46}
\]
Their exact difference is \(1/(24B)\).  The product duals used above match all cross-block observable interactions, so those interactions cannot close the gap. \(\square\)
\end{proof}
\par\smallskip\noindent \textit{Justification.} Tensorization shows that the finite separation survives normalization and unrestricted cross-block competitors at inferential scale; its proof reconstructs the finite primal and observable-moment dual for every degree-restricted space before tensoring the local dual optimizers. \textit{Gap.} The matrix framework in \cite{Middleton2021Unifying} and the optimized quadratic theory in \cite{HarshawMiddletonSavje2026} provide no product law for all-degree observable spaces or exact higher-degree gaps. \textit{Consumer.} Independently randomized small networks or blocks can reuse one certified lookup rule while retaining the exact advantage after aggregation.

\begin{theorem}[thm:studentized-coverage]\label{thm:studentized-coverage}
The bounded-score and nondegeneracy conditions imply \(\sqrt B(\widehat\tau_B-\bar\tau_B)/\sqrt{V_B}\Rightarrow N(0,1)\). Conditional on the denominator law \(\widehat W_B-W_B\xrightarrow{p}0\) and blockwise conservativeness \(b_{Bj}\ge\sigma_{Bj}^2\) for every \(B,j\), one has \(W_B\ge V_B\) and \(\liminf_{B\to\infty}\Pr\{\bar\tau_B\in I_B(\alpha)\}\ge1-\alpha\). If \(W_B\to W<\infty\), coverage converges to \(2\Phi(z_\alpha\sqrt{W/V})-1\); if \(W_B-V_B\to0\), the feasible studentized statistic converges to \(N(0,1)\).
\end{theorem}
\begin{proof}
All expectations, variances, probabilities, and convergence statements below are under the row-specific randomization law. The potential-outcome schedules are fixed, so assignment is the only source of randomness.

For \(1\le j\le B\), define
\[
Y_{Bj}:=X_{Bj}-\tau_{Bj},
\qquad
s_B^2:=\sum_{j=1}^B\sigma_{Bj}^2=B V_B.
\tag{1}
\]
The last equality is the definition of \(V_B\). By the definitions \(\tau_{Bj}=\mathbb E X_{Bj}\) and \(\sigma_{Bj}^2=\operatorname{Var}(X_{Bj})\),
\[
\mathbb E Y_{Bj}=0,
\qquad
\operatorname{Var}(Y_{Bj})=\sigma_{Bj}^2.
\tag{2}
\]
Assumption~\texttt{ass:block-independence} makes \(Y_{B1},\ldots,Y_{BB}\) independent. Assumption~\texttt{ass:uniform-score-bound} and the triangle inequality for expectation give
\[
|\tau_{Bj}|=|\mathbb E X_{Bj}|
\le \mathbb E|X_{Bj}|\le M,
\qquad
|Y_{Bj}|\le |X_{Bj}|+|\tau_{Bj}|\le2M.
\tag{3}
\]
In particular, every \(Y_{Bj}\) has finite variance. Assumption~\texttt{ass:nondegenerate-average-variance} states that \(V_B\to V>0\). Hence, for all sufficiently large \(B\),
\[
V_B\ge V/2,
\qquad
s_B^2=B V_B\ge BV/2>0,
\qquad
s_B\longrightarrow\infty.
\tag{4}
\]
Fix \(\varepsilon>0\). Equations (3)--(4) imply \(2M\le\varepsilon s_B\) for every sufficiently large \(B\). For those rows the events \(\{|Y_{Bj}|>\varepsilon s_B\}\) are empty for every \(j\), and therefore
\[
\frac1{s_B^2}\sum_{j=1}^B
\mathbb E\!\left[Y_{Bj}^2
\mathbf 1\{|Y_{Bj}|>\varepsilon s_B\}\right]=0.
\tag{5}
\]
Thus (5) verifies, rather than assumes, the Lindeberg condition. Lemma~\texttt{lem:lindeberg-feller-independent-array}, applied with row size \(k_B=B\) using (1)--(5), yields
\[
\frac{\sum_{j=1}^B Y_{Bj}}{s_B}\Rightarrow N(0,1).
\tag{6}
\]
Using the definitions of \(\widehat\tau_B\), \(\bar\tau_B\), and \(s_B\), equation (6) is precisely
\[
T_B:=\frac{\sqrt B(\widehat\tau_B-\bar\tau_B)}{\sqrt{V_B}}
=\frac{\sum_{j=1}^B(X_{Bj}-\tau_{Bj})}{\sqrt{B V_B}}
\Rightarrow N(0,1).
\tag{7}
\]

We next derive the denominator law from its primitive row-average variance condition. Let
\[
D_B:=\widehat W_B-W_B
=\frac1B\sum_{j=1}^B(G_{Bj}-b_{Bj}).
\tag{8}
\]
The definition \(b_{Bj}=\mathbb E G_{Bj}\) gives \(\mathbb E D_B=0\). Assumption~\texttt{ass:block-independence} makes \(G_{B1},\ldots,G_{BB}\) independent because they are coordinate projections of independent pairs. Consequently all cross-block covariances in (8) vanish, and
\[
\mathbb E D_B^2
=\operatorname{Var}(D_B)
=\frac1{B^2}\sum_{j=1}^B\operatorname{Var}(G_{Bj})
\longrightarrow0
\tag{9}
\]
by Assumption~\texttt{ass:bound-lln}. For every \(\varepsilon>0\), the pointwise inequality
\[
\varepsilon^2\mathbf 1\{|D_B|>\varepsilon\}\le D_B^2
\tag{10}
\]
and (9) imply
\[
\Pr(|\widehat W_B-W_B|>\varepsilon)
=\Pr(|D_B|>\varepsilon)
\le \frac{\mathbb E D_B^2}{\varepsilon^2}
\longrightarrow0.
\tag{11}
\]
Thus \(\widehat W_B-W_B\xrightarrow{p}0\) is a conclusion. A simple sufficient condition for Assumption~\texttt{ass:bound-lln} is the existence of \(C<\infty\) such that
\[
\sup_{B\ge1}\max_{1\le j\le B}
\mathbb E[(G_{Bj}-b_{Bj})^2]\le C,
\tag{12}
\]
because then the quantity in (9) is at most \(C/B\). In particular, uniformly bounded finite assignment-level lookup tables satisfy (12): if \(|G_{Bj}|\le L<\infty\) uniformly over rows, blocks, and assignment atoms, then \(|b_{Bj}|\le L\) and the left side of (12) is at most \(4L^2\).

Assumption~\texttt{ass:block-conservativeness} and the definitions of \(W_B\) and \(V_B\) give
\[
W_B-V_B
=\frac1B\sum_{j=1}^B(b_{Bj}-\sigma_{Bj}^2)\ge0.
\tag{13}
\]
For all sufficiently large \(B\), define
\[
R_B:=\frac{\max(\widehat W_B,0)}{V_B};
\tag{14}
\]
the division is valid by (4). Fix \(\eta\in(0,1)\). Since \(\max(\widehat W_B,0)\ge\widehat W_B\), equations (4), (11), and (13) imply
\[
\begin{aligned}
\Pr(R_B<1-\eta)
&\le \Pr\!\left(\frac{\widehat W_B}{V_B}<1-\eta\right)\\
&\le \Pr\!\left(|\widehat W_B-W_B|>\eta V_B\right)\\
&\le \Pr\!\left(|\widehat W_B-W_B|>\eta V/2\right)
\longrightarrow0.
\end{aligned}
\tag{15}
\]

Definition~\texttt{def:wald-interval} and equations (7) and (14) give the exact event identity
\[
\{\bar\tau_B\in I_B(\alpha)\}
=\{|T_B|\le z_\alpha\sqrt{R_B}\}.
\tag{16}
\]
Because \(\alpha\in(0,1)\), one has \(z_\alpha=\Phi^{-1}(1-\alpha/2)>0\). On \(\{R_B\ge1-\eta\}\), the event \(\{|T_B|\le z_\alpha\sqrt{1-\eta}\}\) is contained in the event in (16). Hence
\[
\begin{aligned}
\Pr\{\bar\tau_B\in I_B(\alpha)\}
&\ge
\Pr\{|T_B|\le z_\alpha\sqrt{1-\eta}\}
-\Pr(R_B<1-\eta).
\end{aligned}
\tag{17}
\]
By (7), the standard normal law assigns no mass to the two boundary points, so (15)--(17) imply
\[
\liminf_{B\to\infty}\Pr\{\bar\tau_B\in I_B(\alpha)\}
\ge2\Phi\!\left(z_\alpha\sqrt{1-\eta}\right)-1.
\tag{18}
\]
Letting \(\eta\downarrow0\), continuity of \(\Phi\) and the definition of \(z_\alpha\) yield
\[
\liminf_{B\to\infty}\Pr\{\bar\tau_B\in I_B(\alpha)\}
\ge2\Phi(z_\alpha)-1=1-\alpha.
\tag{19}
\]

Suppose next that \(W_B\to W<\infty\). Taking limits in (13) and using \(V_B\to V\) gives \(W\ge V>0\). Equation (11) then implies \(\widehat W_B\xrightarrow{p}W\), so truncation is inactive with probability tending to one and
\[
R_B\xrightarrow{p}r:=W/V\ge1.
\tag{20}
\]
For any \(\delta\in(0,r)\), the probability of \(r-\delta\le R_B\le r+\delta\) tends to one. Therefore
\[
\Pr\{|T_B|\le z_\alpha\sqrt{r-\delta}\}-o(1)
\le \Pr\{\bar\tau_B\in I_B(\alpha)\}
\le \Pr\{|T_B|\le z_\alpha\sqrt{r+\delta}\}+o(1).
\tag{21}
\]
Apply (7) to the two fixed thresholds in (21), and then let \(\delta\downarrow0\). Continuity of \(\Phi\) gives
\[
\Pr\{\bar\tau_B\in I_B(\alpha)\}
\longrightarrow
2\Phi(z_\alpha\sqrt r)-1
=2\Phi\!\left(z_\alpha\sqrt{W/V}\right)-1.
\tag{22}
\]

Finally, suppose \(W_B-V_B\to0\). Equations (4), (11), and this condition give
\[
\frac{\widehat W_B}{V_B}-1
=\frac{\widehat W_B-W_B}{V_B}
+\frac{W_B-V_B}{V_B}
\xrightarrow{p}0,
\qquad
R_B\xrightarrow{p}1.
\tag{23}
\]
In particular, \(\Pr(\widehat W_B>0)\to1\). Define the feasible statistic on that event by
\[
U_B:=\frac{\sqrt B(\widehat\tau_B-\bar\tau_B)}
{\sqrt{\max(\widehat W_B,0)}}
=\frac{T_B}{\sqrt{R_B}},
\tag{24}
\]
and assign it any fixed value when the denominator is zero. Equation (7) makes \(T_B\) tight, while (23) gives \(R_B^{-1/2}-1\xrightarrow{p}0\) on an event whose probability tends to one. Thus, for \(L>0\) and \(\epsilon>0\),
\[
\Pr(|U_B-T_B|>\epsilon)
\le \Pr(|T_B|>L)
+\Pr(|R_B^{-1/2}-1|>\epsilon/L)
+\Pr(R_B=0).
\tag{25}
\]
First choose \(L\) large using tightness and then apply (23); this proves \(U_B-T_B\xrightarrow{p}0\). For every \(x\in\mathbb R\) and \(\epsilon>0\),
\[
\Pr(T_B\le x-\epsilon)-\Pr(|U_B-T_B|>\epsilon)
\le\Pr(U_B\le x)
\le\Pr(T_B\le x+\epsilon)+\Pr(|U_B-T_B|>\epsilon).
\tag{26}
\]
Equations (7), (25)--(26), followed by \(\epsilon\downarrow0\), prove \(U_B\Rightarrow N(0,1)\).

For an exact finite-support stress test showing why no finite-sample coverage is asserted, fix \(\alpha\in(0,1)\). Since \(z_\alpha>0\), choose a rational number \(\rho\) satisfying
\[
0<\rho<\min\{z_\alpha^2,1-\alpha\}.
\tag{27}
\]
In a one-block four-atom randomization, let \(X\) be uniform on \(\{-1,1\}\), let \(H\) be an independent Bernoulli variable with success probability \(\rho\), and put \(G=H/\rho\). Its four atom probabilities are positive and rational. Moreover,
\[
\tau=\mathbb E X=0,
\qquad
\operatorname{Var}(X)=1=\mathbb E G.
\tag{28}
\]
Thus expected conservativeness holds with equality. When \(H=0\), the interval radius is zero and the realized center \(X\in\{-1,1\}\) excludes zero. When \(H=1\), its radius is \(z_\alpha/\sqrt\rho>1\), so it contains zero. Its coverage is therefore exactly \(\rho<1-\alpha\). Under independent replication, \(\operatorname{Var}(G)=\rho^{-1}-1<\infty\), and the primitive denominator quantity is
\[
\frac1{B^2}\sum_{j=1}^B\operatorname{Var}(G_{Bj})
=\frac{\rho^{-1}-1}{B}\longrightarrow0.
\tag{29}
\]
Hence the finite-sample failure is compatible with every asymptotic condition used above and delineates the scope of the Wald conclusion. \(\square\)
\end{proof}
\par\smallskip\noindent \textit{Justification.} Expected conservativeness alone does not imply coverage. This theorem derives the numerator central limit theorem from the primitive independent-block, bounded-score, and nondegeneracy conditions, derives the denominator law from its primitive variance-negligibility condition, and supplies the correct \(B^{-1/2}\) standard-error scaling. \textit{Gap.} The finite-population central limit theory in \cite{LiDing2017FPCLT} is a benchmark and \cite{HarshawMiddletonSavje2026} supplies consistency conditions for quadratic bound estimators, but neither composes an all-degree observable bound with conservative studentization. The declared independent-array Lindeberg--Feller lemma supplies only the classical probability substrate; all verification for the present fixed-potential-outcome design is derived here. \textit{Consumer.} It licenses asymptotic Wald inference for genuinely independent replicated binary exposure blocks using the realized higher-order lookup tables, without asserting finite-sample coverage or a connected-network central limit theorem.

\begin{openendedquestion}[oeq:school-exposure-block]\label{oeq:school-exposure-block}
Can \(\mathcal H_{\mathrm{school}}\) construct a nontrivial finite exposure block from the Paluck--Shepherd--Aronow school-network design whose reduced rational program is lossless for the original assignment law and direct-effect score, whose lifted primal and dual certificates verify the unreduced observable-margin optimum, and whose higher-order optimum strictly improves the best observable quadratic bound without altering the estimand?
\end{openendedquestion}
\par\smallskip\noindent \textit{Justification.} This is the genuine application bottleneck: the finite theory is exact but exponential without exploitable exposure-design structure. \textit{Gap.} The school-data exercise in \cite{HarshawMiddletonSavje2026} is a quadratic empirical illustration based on \cite{PaluckShepherdAronow2016}, not a higher-order reanalysis, and \cite{AronowSamii2017Interference} does not supply an exact all-degree decomposition. \textit{Consumer.} A successful certificate would deliver the first consumer-scale comparison on the school-network substrate while preserving its actual randomization and causal contrast.

\begin{lemma}[lem:lindeberg-feller-independent-array]\label{lem:lindeberg-feller-independent-array}
For every row \(n\), let \(Y_{n1},\ldots,Y_{nk_n}\) be independent real random variables with \(\mathbb E Y_{nj}=0\) and finite variances, and set \(s_n^2=\sum_{j=1}^{k_n}\operatorname{Var}(Y_{nj})>0\). If, for every \(\varepsilon>0\),
\[
 \frac{1}{s_n^2}\sum_{j=1}^{k_n}
 \mathbb E\!\left[Y_{nj}^2\mathbf 1\{|Y_{nj}|>\varepsilon s_n\}\right]
 \longrightarrow 0,
\]
then \(s_n^{-1}\sum_{j=1}^{k_n}Y_{nj}\Rightarrow N(0,1)\).
\end{lemma}
\par\smallskip\noindent \textit{Justification.} This is the classical independent triangular-array central limit theorem used after deriving its Lindeberg premise from the primitive block conditions. \textit{Consumer.} It supplies only the standard probability limit step in the design-based coverage theorem.

\begin{proposition}[prop:witness-hms-sdp-optimum]\label{prop:witness-hms-sdp-optimum}
For \(\mathfrak E_\star\), the direct Harshaw--Middleton--Savje compatible semidefinite program under uniform weighting of \(\Theta\) has the exact value
\[
\min_{\substack{H=H^\top,\ H-A\succeq0\\
H_{03}=H_{04}=H_{14}=H_{23}=H_{24}=H_{45}=0}}
2^{-6}\sum_{\theta\in\Theta}\theta^\top H\theta
=\frac76+\frac{\sqrt{17}}9.
\]
All six diagonal entries are free because every coordinate of \(\mathfrak E_\star\) is observed with positive probability.  Consequently this optimum exceeds \(107/72\) by \((8\sqrt{17}-23)/72\) and exceeds \(55/36\) by \((4\sqrt{17}-13)/36\).
\end{proposition}
\begin{proof}
All calculations are exact in \(\mathbb Q(\sqrt{17})\).  Under the uniform weighting on \(\Theta=\{0,1\}^6\),
\[
\mathbb E_q(\theta_i^2)=\frac12,
\qquad
\mathbb E_q(\theta_i\theta_k)=\frac14\quad(i\ne k).
\]
Therefore, on writing \(R\mathbin\bullet T:=\operatorname{tr}(R^\top T)\),
\[
2^{-6}\sum_{\theta\in\Theta}\theta^\top H\theta
=Q\mathbin\bullet H,
\qquad
Q:=\frac14(I_6+\mathbf 1_6\mathbf 1_6^\top). \tag{4}
\]
Direct substitution of the three score vectors from \(\text{def:witness-experiment}\) into \(\text{def:score-variance}\) gives
\[
9A=
\begin{pmatrix}
8&0&-8&0&-2&2\\
0&6&0&6&-3&-9\\
-8&0&8&0&2&-2\\
0&6&0&6&-3&-9\\
-2&-3&2&-3&2&4\\
2&-9&-2&-9&4&14
\end{pmatrix}. \tag{5}
\]
The jointly unobservable distinct-coordinate pairs are
\[
\mathcal N=\{03,04,14,23,24,45\}. \tag{6}
\]
Every coordinate belongs to at least one displayed observation set, so there is no unobservable diagonal.  For a feasible \(H\), set \(M:=H-A\).  Then \(M\succeq0\), and (5)--(6) together with \(H_{ik}=0\) on \(\mathcal N\) force
\[
(M_{03},M_{04},M_{14},M_{23},M_{24},M_{45})
=\left(0,\frac29,\frac13,0,-\frac29,-\frac49\right). \tag{7}
\]

Put \(\rho:=\sqrt{17}\) and define
\[
S:=
\begin{pmatrix}
\frac12&\frac14&\frac14&0&-\frac\rho{68}&\frac14\\
\frac14&\frac12&\frac14&\frac14&-\frac\rho{34}&\frac14\\
\frac14&\frac14&\frac12&\frac{\rho-1}{8}&\frac{3\rho}{68}&\frac14\\
0&\frac14&\frac{\rho-1}{8}&\frac12&\frac14&\frac14\\
-\frac\rho{68}&-\frac\rho{34}&\frac{3\rho}{68}&\frac14&\frac12&\frac{5\rho}{68}\\
\frac14&\frac14&\frac14&\frac14&\frac{5\rho}{68}&\frac12
\end{pmatrix}. \tag{8}
\]
This matrix has the exact factorization \(S=LDL^\top\), where
\[
L=
\begin{pmatrix}
1&0&0&0&0&0\\
\frac12&1&0&0&0&0\\
\frac12&\frac13&1&0&0&0\\
0&\frac23&\frac{-5+3\rho}{8}&1&0&0\\
-\frac\rho{34}&-\frac\rho{17}&\frac{3\rho}{17}&1+\frac{3\rho}{17}&1&0\\
\frac12&\frac13&\frac14&\frac{3+\rho}{4}&\frac\rho4&1
\end{pmatrix}
\]
and
\[
D=\operatorname{diag}\left(
\frac12,\frac38,\frac13,\frac{5\rho-19}{32},
\frac{9-\rho}{34},0\right). \tag{9}
\]
Multiplication verifies (8)--(9).  Moreover, \(5\rho>19\) because \(25\rho^2=425>361\), and \(\rho<9\) because \(17<81\).  Hence every diagonal entry of \(D\) is nonnegative, so (9) proves \(S\succeq0\).

The matrices \(S\) and \(Q\) agree on every diagonal and on every observable off-diagonal pair.  Thus \(Q-S\) is supported on \(\mathcal N\).  Using (7), with the factor two for symmetric off-diagonal entries, gives
\[
\begin{aligned}
(Q-S)\mathbin\bullet M
={}&2\left[
\left(\frac14+\frac\rho{68}\right)\frac29
+\left(\frac14+\frac\rho{34}\right)\frac13
\right.\\
&\left.\hspace{18mm}
+\left(\frac14-\frac{3\rho}{68}\right)\left(-\frac29\right)
+\left(\frac14-\frac{5\rho}{68}\right)\left(-\frac49\right)
\right]\\
=&\frac{2\rho-1}{18}. \tag{10}
\end{aligned}
\]
The two zero entries in (7) contribute nothing.  The decomposition (9) also shows directly that
\[
S\mathbin\bullet M
=\sum_{\ell=1}^6 D_{\ell\ell}(L^\top M L)_{\ell\ell}
=\sum_{\ell=1}^6D_{\ell\ell}(Le_\ell)^\top M(Le_\ell)\ge0, \tag{11}
\]
where the last inequality uses \(D_{\ell\ell}\ge0\) and \(M\succeq0\).  Finally, (4)--(5) give \(Q\mathbin\bullet A=11/9\).  Combining (10)--(11), every feasible \(H\) obeys
\[
Q\mathbin\bullet H
=Q\mathbin\bullet A+(Q-S)\mathbin\bullet M+S\mathbin\bullet M
\ge \frac{11}{9}+\frac{2\rho-1}{18}
=\frac76+\frac\rho9. \tag{12}
\]

For attainment, set
\[
w:=\begin{pmatrix}1&\frac32&-1&0&\frac\rho2&-2\end{pmatrix}^{\!\top},
\qquad
M^\star:=\frac4{9\rho}ww^\top,
\qquad
H^\star:=A+M^\star. \tag{13}
\]
The positive scalar in (13) makes \(M^\star\succeq0\).  Its entries on \(\mathcal N\) are exactly the six values in (7), so (5) implies that \(H^\star\) satisfies all six zero restrictions.  Thus \(H^\star\) is primal feasible.  Also
\[
\|w\|_2^2=\frac{25}{2},
\qquad
\mathbf 1_6^\top w=\frac{\rho-1}{2},
\]
and hence
\[
Q\mathbin\bullet M^\star
=\frac4{9\rho}\,w^\top Qw
=\frac4{9\rho}\cdot\frac14
\left\{\frac{25}{2}+\frac{(\rho-1)^2}{4}\right\}
=\frac{2\rho-1}{18}. \tag{14}
\]
Equations (12)--(14) prove the exact optimum.  As a complementary check, direct multiplication gives \(Sw=0\), so equality holds in (11).  Subtracting the two benchmark values yields
\[
\left(\frac76+\frac\rho9\right)-\frac{107}{72}
=\frac{8\rho-23}{72},
\qquad
\left(\frac76+\frac\rho9\right)-\frac{55}{36}
=\frac{4\rho-13}{36}. \tag{15}
\]
Both differences are positive: \(8\rho>23\) follows from \(64\rho^2=1088>529\), and \(4\rho>13\) follows from \(16\rho^2=272>169\).
\end{proof}
\par\smallskip\noindent \textit{Justification.} This exact algebraic primal--dual certificate solves the published quadratic semidefinite comparator itself for the six-coordinate witness, rather than inferring only the weaker Boolean quadratic lower bound. \textit{Gap.} The containing Boolean degree-two program has optimum \(55/36\), but strict class inclusion does not determine the optimum of the smaller Harshaw--Middleton--Savje positive-semidefinite class. \textit{Consumer.} Together with the quartic certificate, it gives the exact direct comparator gap \((8\sqrt{17}-23)/72\), approximately \(0.138678403\).

\section*{Technical internal limitation (diagnostic only)}

Exact enumeration uses \(2^K\) domination inequalities, \(|\mathcal J|\) coefficients, and assignment truth tables of total size \(\sum_z2^{|O_z|}\). Independent product designs tensorize exactly, but this gives no lossless decomposition for a large connected school network. Nonnegative expected bounds also need not admit nonnegative assignmentwise tables, although the displayed quartic witness does. Finiteness of every row does not by itself give a triangular-array denominator law: the inference result requires \(B^{-2}\sum_j\operatorname{Var}(G_{Bj})\to0\); uniformly bounded lookup tables are a simple sufficient condition, but rare observation probabilities can destroy uniform boundedness. The numerator proof is specific to independent product rows; it does not establish a central limit theorem for connected dependent-exposure arrays. If \(q\) is not complement invariant, averaging a bound with its complemented version can change the objective; the odd-degree collapse then supplies no degree-three limitation, and no unsymmetric cubic-equals-quadratic claim is made.

\section{Honest scope}

The exact \(107/72\), \(55/36\), and \(7/6+\sqrt{17}/9\) witness values have finite or algebraic certificates; the school-network exposure-block question remains open. The complete observable-function and Pareto results concern fixed binary schedules and finite randomized designs with arbitrary real probabilities on their support; rationality is required only for exact rational optimizer and assignment-table certificates and for the displayed witness. The binary Aronow--Samii specialization records a published conservative Horvitz--Thompson baseline and is not claimed as a new correction. There is no finite-sample coverage claim, no universal interval-shortening claim, no efficient general large-design algorithm, and no extension to real outcomes. Conservative studentized coverage is an asymptotic independent-block result. Its numerator central limit theorem follows from the declared independent-array Lindeberg--Feller lemma after block independence, bounded scores, and \(V_B\to V>0\) establish all of that lemma's premises, including the Lindeberg condition. Its denominator weak law follows separately from \(B^{-2}\sum_j\operatorname{Var}(G_{Bj})\to0\); uniformly bounded finite lookup tables are sufficient but are not automatic across a triangular array. A four-atom finite experiment satisfies all replicated asymptotic primitives while failing nominal coverage at \(B=1\), so no finite-sample inference conclusion is implicit. The direct Harshaw--Middleton--Savje calculation is an exact six-dimensional witness semidefinite program, not a consumer-scale computation. Under full joint observability the theory reduces to exact variance estimability; under the stated complement symmetries and observation sets of size at most three it reduces to quadratic optimality, while no odd-degree collapse is claimed for asymmetric objective weights. The Harshaw--Middleton--Savje school-data analysis is only a quadratic empirical illustration, not a reanalysis of the original causal conclusions.

\begin{thebibliography}{99}

  \bibitem{HarshawMiddletonSavje2026} Christopher Harshaw, Joel A. Middleton, and Fredrik S\"avje (2026), Optimized Variance Estimation under Interference and Complex Experimental Designs, Journal of the American Statistical Association, DOI: 10.1080/01621459.2026.2627027; arXiv:2112.01709.

  \bibitem{Middleton2021Unifying} Joel A. Middleton (2021), Unifying Design-based Inference: On Bounding and Estimating the Variance of any Linear Estimator in any Experimental Design, arXiv:2109.09220.

  \bibitem{Lu2019FactorialBinary} Jiannan Lu (2019), Sharpening Randomization-Based Causal Inference for \(2^2\) Factorial Designs with Binary Outcomes, Statistical Methods in Medical Research 28(4), 1064--1078, DOI: 10.1177/0962280217745720; arXiv:1711.04432.

  \bibitem{AronowGreenLee2014} Peter M. Aronow, Donald P. Green, and Donald K. K. Lee (2014), Sharp Bounds on the Variance in Randomized Experiments, Annals of Statistics 42, 850--871, DOI: 10.1214/13-AOS1200; arXiv:1405.6555.

  \bibitem{AronowSamii2013} Peter M. Aronow and Cyrus Samii (2013), Conservative Variance Estimation for Sampling Designs with Zero Pairwise Inclusion Probabilities, Survey Methodology 39(1), 231--241.

  \bibitem{AronowSamii2017Interference} Peter M. Aronow and Cyrus Samii (2017), Estimating Average Causal Effects under General Interference, with Application to a Social Network Experiment, Annals of Applied Statistics 11, 1912--1947; arXiv:1305.6156.

  \bibitem{PaluckShepherdAronow2016} Elizabeth Levy Paluck, Hana Shepherd, and Peter M. Aronow (2016), Changing Climates of Conflict: A Social Network Experiment in 56 Schools, Proceedings of the National Academy of Sciences 113, 566--571.

  \bibitem{LiDing2017FPCLT} Xinran Li and Peng Ding (2017), General Forms of Finite Population Central Limit Theorems with Applications to Causal Inference, Journal of the American Statistical Association 112, 1759--1769; arXiv:1610.04821.

  \bibitem{ChattopadhyayImbens2025} Ambarish Chattopadhyay and Guido W. Imbens (2025), Neymanian Inference in Randomized Experiments, arXiv:2409.12498.

  \bibitem{AzrielKriegerKapelner2026} David Azriel, Abba M. Krieger, and Adam Kapelner (2026), Optimal Designs with Robust Inference for Binary Treatment Effects, arXiv:2607.05768.

  \bibitem{Billingsley1995ProbabilityMeasure} Patrick Billingsley (1995), Probability and Measure, third edition, Wiley, Theorem 27.2, DOI: 10.1002/0470316962.

\end{thebibliography}

\end{document}